\documentclass[12pt,oneside]{memoir}

\usepackage[latinica,biblatex]{matfmaster}

% Više prostora između zaglavlja (header) i teksta
\setheaderspaces{*}{1.7\onelineskip}{*}
\checkandfixthelayout

\usepackage{iftex}
\ifpdftex
  \usepackage{cmsrb}
\fi
\usepackage{listings}
\usepackage{amsmath}
\usepackage{amssymb}
\usepackage{booktabs}
\usepackage{multirow}
\usepackage{xcolor}

\lstset{
  basicstyle=\ttfamily\small,
  keywordstyle=\color{blue}\bfseries,
  commentstyle=\color{gray},
  stringstyle=\color{red!70!black},
  showstringspaces=false,
  breaklines=true,
  frame=single,
  numbers=left,
  numberstyle=\tiny\color{gray},
  tabsize=4,
}

% Prelomljiv tt font za duga imena fajlova (prelama na _ . /)
\DeclareUrlCommand\code{\urlstyle{tt}}

\bib{docs/master}

\autor{Dušan Trtica}
\naslov{Primena tehnika zadovoljenja ograničenja i celobrojnog programiranja na problem generisanja rasporeda časova}
\godina{2026}
\mentor{dr Filip \textsc{Marić}, redovni profesor\\ Univerzitet u Beogradu, Matematički fakultet}
\komisijaA{dr Aleksandar \textsc{Kartelj}, vanredni profesor\\ Univerzitet u Beogradu, Matematički fakultet}
\komisijaB{dr Miljan \textsc{Knežević}, docent\\ Univerzitet u Beogradu, Matematički fakultet}

\apstr{%
U ovom radu razmatramo problem automatskog generisanja rasporeda časova
na univerzitetskom nivou, na podacima Matematičkog fakulteta
Univerziteta u Beogradu. Rad je sproveden u tri koraka. Najpre su
razvijene dve formulacije osnovnog problema, u kojima se traži raspored
koji zadovoljava tvrda ograničenja: jedna u paradigmi programiranja
ograničenja (CP-SAT) i jedna kao binarni celobrojni linearni model
(0-1 ILP), poseban slučaj mešovitog celobrojnog programiranja, rešavan
SCIP-om (oznaka MIP/SCIP); obe su implementirane pomoću biblioteke
Google OR-Tools. Zatim je sprovedeno njihovo eksperimentalno poređenje
na tri skale realnih podataka, po vremenu konstrukcije modela i
rešavanja, veličini modela i memorijskom otisku. Obe formulacije daju
validna rešenja na svim skalama, pri čemu CP-SAT do dopustivog
rasporeda dolazi brže, te je iz praktičnih razloga izabran kao
osnova za nastavak rada. U trećem koraku CP model je proširen:
uvedena je arhitektura pravila sa tvrdim i mekim ograničenjima i
funkcijom cilja, dodata su ograničenja kvaliteta rasporeda -- spajanje
časova u blokove, jedna lokacija po danu, izbegavanje procepa i
kapacitet učionica -- i raspoređivanje nastavnog osoblja. Prošireni
model je ponovo evaluiran, uz praćenje vrednosti funkcije cilja kroz
vreme i dodatno merenje u okruženju sa manjim brojem
termina.
}

\kljucnereci{raspored časova, programiranje ograničenja, celobrojno programiranje,
CP-SAT, MIP, OR-Tools, optimizacija, Matematički fakultet}

\begin{document}
% ==============================================================================
\frontmatter
% ==============================================================================
\naslovna
\komisija
\posveta{Mojoj porodici}
\apstrakt
\tableofcontents*

% ==============================================================================
\mainmatter
% ==============================================================================

% ------------------------------------------------------------------------------
\chapter{Uvod}
\label{chp:uvod}
% ------------------------------------------------------------------------------

Generisanje rasporeda časova predstavlja jedan od najčešćih i najpraktičnijih
kombinatornih problema sa kojim se suočavaju obrazovne institucije širom sveta.
Svake godine, na početku akademskog semestra, univerziteti i fakulteti moraju
da rasporede stotine predavanja, vežbi i laboratorijskih časova u ograničen
broj učionica, uzimajući u obzir radne dane, raspoložive termine i različite
vrste ograničenja.

Problem raspoređivanja časova (eng. \emph{University Course Timetabling
Problem} -- UCTP) jeste NP-težak problem \cite{even1975, garey1979}, što
znači da ne postoji poznat algoritam koji bi ga efikasno rešio u opštem
slučaju. Tradicionalno, rasporedi se prave ručno, što zahteva veliku
količinu ljudskog rada, vremena i stručnosti. Čak i iskusni raspoređivači
mogu da propuste kolizije ili da stvore neoptimalan raspored.

U ovom radu istražujemo dva pristupa za automatsko rešavanje UCTP-a:

\begin{enumerate}
\item \textbf{Programiranje ograničenja} (eng. \emph{Constraint Programming}
  -- CP) -- deklarativna paradigma u kojoj se problem definiše preko skupa
  promenljivih, njihovih domena i ograničenja koja moraju biti zadovoljena.
  Koristimo Google OR-Tools CP-SAT rešavač \cite{ortools, cpsat2020}.

\item \textbf{Mešovito celobrojno programiranje} (eng. \emph{Mixed Integer
  Programming} -- MIP) -- matematička optimizacija u kojoj se problem
  zapisuje linearnom funkcijom cilja i linearnim ograničenjima nad
  promenljivama koje mogu biti celobrojne ili kontinualne. Model koji
  razvijamo u ovom radu koristi isključivo binarne promenljive, pa je
  reč o binarnom celobrojnom linearnom modelu (0-1 ILP, odnosno BIP),
  posebnom slučaju šire klase MIP modela. Rešavamo ga SCIP
  rešavačem \cite{scip2009} integrisanim u OR-Tools.
\end{enumerate}

\section{Cilj rada}

Rad ima tri cilja koja se nadovezuju jedan na drugi.

\begin{enumerate}
\item \textbf{Dve formulacije osnovnog problema.} Za isti skup tvrdih
  ograničenja razvijamo CP i MIP verziju, svaku u idiomu njene
  paradigme.

\item \textbf{Eksperimentalno poređenje.} Obe formulacije poredimo na
  tri skale realnih podataka (mali, srednji i veliki podskup) po
  vremenu konstrukcije modela, vremenu rešavanja, memorijskom otisku,
  veličini modela i validnosti dobijenog rasporeda.

\item \textbf{Proširenje i ponovna evaluacija.} Jednu od dve
  formulacije proširujemo mekim ograničenjima, funkcijom cilja i
  raspoređivanjem nastavnog osoblja, pa prošireni model ponovo
  evaluiramo.
\end{enumerate}

Veličinu modela (broj promenljivih i ograničenja) posmatramo kao
opisnu karakteristiku formulacije, a ne kao meru njenog kvaliteta:
dve paradigme problem zapisuju na različite načine, pa
sam broj promenljivih nije uporediv između njih. Kao merila koristimo
ukupno vreme do dopustivog rešenja, ponašanje pri rastu instance i
nezavisno proverenu validnost rasporeda.

Za nastavak rada izabran je CP model. Izbor je pragmatičan i vezan za
konkretan problem, instance i implementacije opisane u ovom radu:
CP-SAT je na svim ispitanim skalama brže dolazio do dopustivog
rasporeda, a njegov zapis se pokazao pogodnim za postupno dodavanje
novih pravila. MIP formulacija takođe daje validne rasporede,
a drugačija formulacija ili drugi rešavač mogli bi dati drugačiji
odnos rezultata.

\section{Struktura rada}

Rad je organizovan na sledeći način:

\begin{itemize}
\item U glavi \ref{chp:teorija} izlažemo teorijske osnove: definiciju
  problema raspoređivanja časova, paradigmu programiranja ograničenja
  i paradigmu celobrojnog linearnog programiranja.

\item U glavi \ref{chp:alati} dajemo kratak pregled korišćenih alata:
  biblioteke Google OR-Tools, sistema za gradnju Bazel i pomoćnih
  Python biblioteka.

\item U glavi \ref{chp:formulacija} izlažemo model ulaznih podataka,
  formalnu specifikaciju problema u obe paradigme -- promenljive,
  domene i ograničenja -- i njihovu implementaciju.

\item U glavi \ref{chp:evaluacija} predstavljamo preliminarnu evaluaciju
  osnovnog modela: metodologiju merenja, poređenje dve formulacije na tri
  skale problema i analizu dobijenih rezultata. Ova evaluacija služi i
  kao osnov za izbor formulacije koju dalje razvijamo.

\item U glavi \ref{chp:nadogradnja} CP model, izabran u prethodnoj glavi,
  nadograđujemo arhitekturom pravila, funkcijom cilja i uvodimo nova ograničenja.

\item U glavi \ref{chp:osoblje} model proširujemo nastavnim osobljem:
  dodelom nastave, zajedničkim sesijama i pravilima koja štite
  raspored nastavnika.

\item U glavi \ref{chp:finalna-evaluacija} dajemo završnu evaluaciju
  proširenog modela na tri skale, sa naglaskom na vrednost funkcije
  cilja u zavisnosti od proteklog vremena izračunavanja.

\item U glavi \ref{chp:zakljucak} sumiramo zaključke i navodimo pravce
  daljeg rada.
\end{itemize}


% ------------------------------------------------------------------------------
\chapter{Teorijske osnove}
\label{chp:teorija}
% ------------------------------------------------------------------------------

U ovoj glavi izlažemo teorijski okvir za razumevanje problema i pristupa
koje koristimo u ovom radu. Najpre definišemo problem raspoređivanja
časova, zatim opisujemo paradigmu programiranja ograničenja, i na kraju
predstavljamo celobrojno linearno programiranje.

\section{Problem raspoređivanja časova}
\label{sec:uctp}

Problem raspoređivanja univerzitetskih časova (UCTP) sastoji se od
dodele časova predavanja, vežbi u raspoložive
vremenske termine i učionice, tako da budu zadovoljena određena
ograničenja \cite{schaerf1999, burke2004}.

Formalno, problem se može opisati sledećim elementima:

\begin{itemize}
\item \textbf{Skup sesija} $\mathcal{S} = \{s_1, s_2, \ldots, s_n\}$
  -- časovi koje treba rasporediti. Svaka sesija pripada određenom
  predmetu i grupi studenata.

\item \textbf{Skup vremenskih termina} $\mathcal{T} = \{t_1, t_2, \ldots, t_m\}$
  -- raspoloživi termini u toku radne nedelje, određeni kombinacijom
  radnog dana i sata.

\item \textbf{Skup učionica} $\mathcal{R} = \{r_1, r_2, \ldots, r_k\}$
  -- prostorije u kojima se nastava može održati.

\item \textbf{Skup ograničenja} $\mathcal{C}$ -- pravila koja raspored
  mora da zadovolji.
\end{itemize}

Cilj je pronaći funkciju dodele $\sigma : \mathcal{S} \to \mathcal{T}
\times \mathcal{R}$ takvu da su sva ograničenja iz $\mathcal{C}$
zadovoljena.

\subsection{Vrste ograničenja}
\label{subsec:ogranicenja}

Ograničenja se u literaturi klasifikuju u dve kategorije
\cite{schaerf1999, kingston2013}:

\begin{enumerate}
\item \textbf{Tvrda ograničenja} (eng. \emph{hard constraints}) --
  ograničenja koja moraju biti bezuslovno zadovoljena. Rešenje koje
  krši neko tvrdo ograničenje nije dopustivo. Primeri:
  \begin{itemize}
  \item Nikoje dve sesije ne mogu se održavati u istoj učionici u
    isto vreme.
  \item Jedna grupa studenata ne može istovremeno pohađati dva
    različita predavanja.
  \item Sesija koja zahteva računare mora biti raspoređena u
    učionicu opremljenu računarima.
  \end{itemize}

\item \textbf{Meka ograničenja} (eng. \emph{soft constraints}) --
  ograničenja čije zadovoljenje je poželjno, ali ne i obavezno.
  Njihovo kršenje ne čini rešenje nedopustivim, već samo lošijim.
  Primeri:
  \begin{itemize}
  \item Ravnomerna raspodela časova tokom nedelje.
  \item Minimizacija broja promena lokacija za jednu grupu u toku dana.
  \item Izbegavanje procepa (eng. \emph{gap}) -- slobodnih sati između
    dva časa -- u rasporedu za studente.
  \end{itemize}
\end{enumerate}

U prvom delu ovog rada bavimo se isključivo tvrdim ograničenjima, tj.
pronalaženjem dopustivog rasporeda (eng. \emph{feasibility problem}).
Proširenje mekim ograničenjima i optimizacija kvaliteta rasporeda
predviđeni su za drugi deo rada.

\subsection{NP-težina problema}

Even, Itai i Shamir \cite{even1975} su dokazali da je problem
raspoređivanja NP-težak redukcijom sa problema bojenja grafova.
Garey i Johnson \cite{garey1979} su ovaj rezultat uključili u svoju
obimnu katalogizaciju NP-kompletnih problema. Ovo znači da je
malo verovatno da postoji algoritam polinomijalne složenosti za
rešavanje opšteg slučaja problema raspoređivanja, te se u praksi
pribegava heuristikama ili egzaktnim metodama sa eksponencijalnom
najgorom složenošću.

\section{Programiranje ograničenja}
\label{sec:cp}

Programiranje ograničenja (eng. \emph{Constraint Programming} -- CP)
je deklarativna paradigma za rešavanje kombinatornih problema
\cite{rossi2006, apt2003}. Problem se formuliše kao \emph{problem
zadovoljenja ograničenja} (eng. \emph{Constraint Satisfaction
Problem} -- CSP).

\subsection{Problem zadovoljenja ograničenja (CSP)}

CSP se formalno definiše trojkom $(\mathcal{X}, \mathcal{D},
\mathcal{C})$ gde je:

\begin{itemize}
\item $\mathcal{X} = \{x_1, x_2, \ldots, x_n\}$ -- skup promenljivih,
\item $\mathcal{D} = \{D_1, D_2, \ldots, D_n\}$ -- skup domena, gde
  $D_i$ sadrži moguće vrednosti promenljive $x_i$,
\item $\mathcal{C} = \{C_1, C_2, \ldots, C_m\}$ -- skup ograničenja,
  gde svako ograničenje $C_j$ specificira dopustive kombinacije
  vrednosti nad podskupom promenljivih.
\end{itemize}

Rešenje CSP-a je dodela vrednosti svim promenljivama iz njihovih
domena tako da su sva ograničenja zadovoljena.

\subsection{Tehnike rešavanja}

CP rešavači kombinuju sistematsku pretragu sa tehnikama propagacije
ograničenja \cite{rossi2006}:

\begin{enumerate}
\item \textbf{Vraćanje unazad} (eng. \emph{backtracking}) --
  pretraga prostora rešenja dodeljujući vrednosti promenljivama
  jednu po jednu. Kada se naiđe na ograničenje koje ne može biti
  zadovoljeno, vrši se vraćanje na poslednju tačku izbora.

\item \textbf{Propagacija ograničenja} (eng. \emph{constraint propagation})
  -- smanjivanje domena promenljivih na osnovu ograničenja, pre i
  tokom pretrage. Ovo eliminiše nedopustive vrednosti rano i
  dramatično smanjuje prostor pretrage.

\item \textbf{Konzistentnost lukova} (eng. \emph{arc consistency}) --
  tehnika propagacije koja obezbeđuje da za svaku vrednost u domenu
  jedne promenljive postoji kompatibilna vrednost u domenu svake
  povezane promenljive.
\end{enumerate}

\subsection{Globalna ograničenja}

Ključna prednost CP-a je podrška za \emph{globalna ograničenja} --
specijalizovane strukture koje obuhvataju odnose između većeg broja
promenljivih i koje imaju efikasne algoritme propagacije \cite{regin1994}.

\textbf{AllDifferent.}
Ograničenje $\texttt{AllDifferent}(x_1, x_2, \ldots, x_n)$ zahteva
da sve promenljive $x_1, \ldots, x_n$ imaju međusobno različite
vrednosti. Dok bi se isto moglo izraziti sa $\binom{n}{2}$ parnih
nejednakosti, globalno ograničenje pruža značajno jaču propagaciju
zahvaljujući algoritmima baziranim na sparivanju u bipartitnim
grafovima \cite{regin1994}.

\textbf{AllowedAssignments.}
Ograničenje $\texttt{AllowedAssignments}$ (poznato i kao
\emph{table constraint}) eksplicitno navodi dozvoljene kombinacije
vrednosti promenljivih pomoću tabele (skupa torki). Ovo je korisno
za ograničenja koja se ne mogu izraziti prostom aritmetikom, na primer
ograničenje da sesija može biti raspoređena samo u podskup učionica.

\subsection{CP-SAT pristup}

Moderna generacija CP rešavača, kakav je Google CP-SAT \cite{cpsat2020},
kombinuje tehnike programiranja ograničenja sa tehnologijama SAT
rešavača. Ovaj pristup se zasniva na \emph{lenjom generisanju klauza}
(eng. \emph{lazy clause generation}) \cite{ohrimenko2009}, gde se
ograničenja propagiraju kao u klasičnom CP, ali se konflikti pamte
u obliku SAT klauza, što sprečava ponavljanje istih pogrešnih izbora
tokom pretrage. Ovo kombinuje snagu globalne propagacije sa efikasnošću
učenja konflikata poznatom iz SAT rešavača.

\section{Celobrojno linearno programiranje}
\label{sec:mip}

Mešovito celobrojno programiranje (eng. \emph{Mixed Integer Programming}
-- MIP) je paradigma matematičke optimizacije u kojoj se problem
formuliše pomoću linearne funkcije cilja i linearnih ograničenja
nad promenljivama koje mogu biti celobrojne ili kontinualne
\cite{wolsey1998, nemhauser1988}.

\subsection{Definicija MIP problema}

MIP problem u opštem obliku glasi:

\begin{align}
\min \quad & \mathbf{c}^T \mathbf{x} \label{eq:mip_obj} \\
\text{p.o.} \quad & A\mathbf{x} \leq \mathbf{b} \label{eq:mip_constr} \\
& x_i \in \mathbb{Z}, \quad i \in I \label{eq:mip_int} \\
& x_j \geq 0, \quad j \notin I \label{eq:mip_cont}
\end{align}

gde je $\mathbf{c}$ vektor koeficijenata funkcije cilja, $A$ matrica
koeficijenata ograničenja, $\mathbf{b}$ vektor desnih strana, a $I$
skup indeksa celobrojnih promenljivih.

Unutar ove klase razlikujemo uže slučajeve. Ako su sve promenljive
celobrojne, odnosno ako je $I$ skup svih indeksa, reč je o
\emph{celobrojnom linearnom programu} (eng. \emph{Integer Linear
Program} -- ILP). Ako su uz to sve promenljive binarne, $x_i \in \{0,
1\}$, govori se o \emph{binarnom celobrojnom linearnom programu}
(0-1 ILP, odnosno BIP).

Formulacija razvijena u ovom radu (odeljak \ref{sec:mip_form}) pripada
poslednjoj klasi: sve njene promenljive su binarne, pa je reč o 0-1
ILP modelu, posebnom slučaju opšteg MIP oblika
\eqref{eq:mip_obj}--\eqref{eq:mip_cont}. U tabelama i u tekstu zadržana
je kraća oznaka MIP/SCIP, koja imenuje paradigmu i rešavač.

Osnovni model, izložen u glavi \ref{chp:formulacija}, postavljen je kao
problem dopustivosti: funkcija cilja iz izraza \eqref{eq:mip_obj} se ne
koristi, već se traži bilo koja dodela vrednosti koja zadovoljava sva
linearna ograničenja. Prošireni CP model iz glava
\ref{chp:nadogradnja} i \ref{chp:osoblje}, nasuprot tome, uvodi meka
ograničenja i funkciju cilja \eqref{eq:objective} koja se minimizuje.

\subsection{Tehnike rešavanja}

MIP rešavači koriste kombinaciju sledećih tehnika \cite{wolsey1998}:

\begin{enumerate}
\item \textbf{LP relaksacija} -- opuštanje celobrojnih ograničenja
  i rešavanje kontinualnog linearnog programa. Ovo daje donju granicu
  za funkciju cilja (ili test nedopustivosti).

\item \textbf{Grananje i odsecanje} (eng. \emph{branch and bound}) --
  sistematska pretraga koja deli prostor rešenja na manje podprobleme,
  koristeći LP relaksaciju za odsecanje grana koje ne mogu sadržati
  bolje rešenje.

\item \textbf{Ravni odsecanja} (eng. \emph{cutting planes}) --
  dodavanje linearnih ograničenja koja eliminišu frakciona rešenja
  LP relaksacije, a ne eliminišu celobrojna rešenja.
\end{enumerate}

\subsection{SCIP rešavač}

SCIP (Solving Constraint Integer Programs) \cite{scip2009} je jedan
od najsnažnijih akademskih rešavača za MIP i mešovito celobrojno
nelinearno programiranje. SCIP kombinuje tehnike grananja i odsecanja
sa propagacijom ograničenja i heuristikama. U našem radu koristimo
SCIP integrisanog u Google OR-Tools putem interfejsa
\texttt{pywraplp}.

Time smo postavili pojmovni okvir: problem raspoređivanja i dva načina
da se on zapiše, jedan preko promenljivih sa domenima i ograničenja
koja se propagiraju, drugi preko binarnih promenljivih i linearnih
nejednakosti. Naredna glava ukratko predstavlja alate kojima su ta dva
zapisa implementirana.


% ------------------------------------------------------------------------------
\chapter{Korišćeni alati i tehnologije}
\label{chp:alati}
% ------------------------------------------------------------------------------

U ovoj glavi dajemo kratak pregled alata i tehnologija korišćenih u
implementaciji, u meri u kojoj su potrebni za razumevanje narednih
glava. Projekat je napisan u programskom jeziku Python \cite{python} i
koristi Bazel \cite{bazel} kao sistem za gradnju i upravljanje
zavisnostima.

\section{Google OR-Tools}
\label{sec:ortools}

Google OR-Tools \cite{ortools} je biblioteka otvorenog koda za
kombinatornu optimizaciju razvijena od strane Google-a. Nudi rešavače
za više tipova problema, uključujući:

\begin{itemize}
\item \textbf{CP-SAT} -- rešavač za programiranje ograničenja sa
  SAT tehnikama, opisan u odeljku \ref{sec:cp}. U Pythonu se koristi
  putem modula \texttt{ortools.sat.python.cp\_model}. Ključne klase
  su \texttt{CpModel} (za definisanje promenljivih i ograničenja)
  i \texttt{CpSolver} (za pokretanje pretrage).

\item \textbf{Linear solver wrapper (pywraplp)} -- uniformni interfejs
  za više MIP rešavača, od kojih u ovom radu koristimo SCIP. Ključna
  klasa je \texttt{Solver}, sa metodama za kreiranje binarnih
  promenljivih (\texttt{BoolVar}), dodavanje linearnih ograničenja
  (\texttt{Add}) i pokretanje rešavanja (\texttt{Solve}).
\end{itemize}

OR-Tools je izabran zato što je otvorenog izvornog koda i što kroz
jedinstven Python API pokriva obe paradigme, pa se dve formulacije
mogu porediti bez razlika koje bi poticale od različitih biblioteka.

\section{Bazel}
\label{sec:bazel}

Bazel \cite{bazel} je sistem za gradnju i testiranje softvera. U ovom
projektu koristi se zbog hermetičkih, reproducibilnih gradnji: iste
verzije zavisnosti i isti izvorni kod daju isti rezultat na svakoj
mašini, što je važno da bi merenja iz glava \ref{chp:evaluacija} i
\ref{chp:finalna-evaluacija} bila ponovljiva. Inkrementalne gradnje
pritom skraćuju ciklus razvoja.

\section{Python biblioteke}
\label{sec:pylibs}

Ulazni podaci se učitavaju i validiraju pomoću biblioteke Pydantic
\cite{pydantic}, koja Python anotacije tipova koristi kao specifikaciju
ulaza. Svaka entitetska klasa (\texttt{Settings}, \texttt{Classroom},
\texttt{Course}, itd.) definisana je kao \texttt{pydantic.dataclass},
pa neispravan ulaz otkazuje već pri parsiranju, a ne u rešavaču;
deklaracija \texttt{Field(alias=...)} pri tome preslikava
\texttt{camelCase} imena iz JSON-a u \texttt{snake\_case} imena u
Python kodu. Konkretan model ulaznih podataka daje odeljak
\ref{sec:input_model}.

Pored toga, openpyxl \cite{openpyxl} koristi se za izvoz rasporeda u
Excel fajl sa posebnim radnim listom za svaku grupu studenata,
PyFunctional \cite{pyfunctional} za fluent filtriranje listi predmeta i
učionica, a pytest \cite{pytest} za testove koji pokrivaju parsiranje
ulaza, formiranje grupa i sesija i proveru da rešavači daju dopustiva
rešenja. Testovi se pokreću kroz Bazel.

Teorijski okvir iz glave \ref{chp:teorija} i alati opisani ovde daju
sve što je potrebno da se problem raspoređivanja zapiše na dva načina.
U narednoj glavi prelazimo sa opšteg opisa na konkretan model: najpre
na ulazne podatke i sesije, zatim na formalnu specifikaciju CP i MIP
modela, i na kraju na njihovu implementaciju u OR-Tools-u.


% ------------------------------------------------------------------------------
\chapter{Formulacija problema i implementacija}
\label{chp:formulacija}
% ------------------------------------------------------------------------------

U ovoj glavi prelazimo sa opšteg opisa problema na konkretan model.
Najpre opisujemo ulazne podatke i način na koji se od njih formiraju
sesije koje treba rasporediti. Zatim dajemo formalnu specifikaciju --
promenljive, njihove domene i ograničenja -- odvojeno za CP i za MIP
paradigmu, i pokazujemo da dve formulacije opisuju isti skup dopustivih
rešenja. Tek nakon toga prikazujemo implementaciju, odnosno kod kojim
se te specifikacije prevode u modele dva rešavača.

\section{Ulazni podaci i sesije}

\subsection{Model ulaznih podataka}
\label{sec:input_model}

Ulazni podaci su definisani u JSON formatu i modelovani pomoću
Pydantic klasa u modulu \code{model.py}; njihov zapis u kodu dat je
uz ostatak implementacije, u listingu \ref{lst:pydantic_model}.

\medskip
Struktura ulaznih podataka je sledeća:

\begin{enumerate}
\item \textbf{Settings} -- podešavanja rasporeda:
  \begin{itemize}
  \item \texttt{workingDays} -- lista radnih dana
    (Ponedeljak, Utorak, ..., Petak),
  \item \texttt{startHour}, \texttt{endHour} -- početak i kraj
    radnog vremena (8--20),
  \item \texttt{duration} -- trajanje jednog časa predavanja i vežbi (1 sat).
  \end{itemize}

\item \textbf{Location} -- lokacije/zgrade fakulteta (Studentski trg,
  Svetog Nikole, Jagićeva).

\item \textbf{Classroom} -- učionice sa atributima:
  \begin{itemize}
  \item identifikator, naziv, lokacija,
  \item da li poseduje računare (\texttt{hasComputers}),
  \item kapacitet.
  \end{itemize}

\item \textbf{Smer} -- smerovi studija (Profesor matematike i
  računarstva, Teorijska matematika, Informatika, itd.). U izvornom
  kodu odgovara klasi \texttt{StudyTrack}, odnosno polju
  \texttt{trackId} u JSON ulazu.

\item \textbf{Course} -- predmeti sa:
  \begin{itemize}
  \item identifikatorom, nazivom, semestrom, smerom,
  \item kvotom (\texttt{Quota}) -- koliko časova predavanja i vežbi
    nedeljno,
  \item zahtevom za računarima (\texttt{needsComputers}), zadatim
    \emph{odvojeno za predavanja i za vežbe}.
  \end{itemize}

  Predmet je pri tome samo nosilac zajedničkog identiteta i kvote:
  predavanja i vežbe se dalje raspoređuju kao nezavisni časovi, o čemu
  detaljno govori sekcija \ref{sec:sesije}.

\item \textbf{StudentsEnrolled} -- broj upisanih studenata po smeru
  i semestru.
\end{enumerate}

Sve navedene celine objedinjuje korenska klasa
\texttt{SchedulingInput}, koja odgovara jednom ulaznom JSON fajlu.
Polja \texttt{rules} i \texttt{groups} imaju podrazumevano prazne
vrednosti i uvode se u kasnijim glavama: \texttt{rules} nosi
konfiguraciju mekih i tvrdih pravila (glava \ref{chp:nadogradnja}), a
\texttt{groups} omogućava da se grupe studenata zadaju eksplicitno
umesto da se izvode iz broja upisanih; u realnom ulazu grupe dolaze iz
fajla sa nastavnim osobljem (odeljak \ref{sec:grupe_iz_osoblja}), koji
je modelovan zasebnom klasom \texttt{StaffInput} (odeljak
\ref{sec:staff_model}). Alias \texttt{departments} zadržan je radi
kompatibilnosti sa starijim ulaznim fajlovima.

\subsection{Realni podaci i poreklo ulaza}
\label{sec:poreklo_podataka}

Sve evaluacije u ovom radu koriste jedan skup podataka, koji
predstavlja letnji semestar nastave na Matematičkom fakultetu
(\code{input_full_2_semester.json}, uz \code{staff_2_semester.json}):

\begin{itemize}
\item 6 smerova,
\item 134 predmeta raspoređena po semestrima (2, 4, 6 i 8),
\item 29 učionica na 3 lokacije,
\item 919 upisanih studenata, iz kojih se dobija 30 grupa deljenjem po
  veličini, odnosno 44 grupe kada se imena preuzmu iz dodela nastave,
\item 112 nastavnika i 264 dodele nastave,
\item radna nedelja od 5 dana, po 12 sati dnevno (8:00--20:00).
\end{itemize}

Za ocenu rezultata važno je razdvojiti šta je preuzeto iz zvaničnih
izvora, a šta je procena. Iz plana i programa studija preuzeti su
smerovi, predmeti, semestri u kojima se predmeti slušaju i nedeljne
kvote časova predavanja i vežbi. Iz spiska prostorija preuzeti su
učionice, njihove lokacije i podatak o tome koje su opremljene
računarima.

Procene su, sa druge strane, kapaciteti učionica, broj upisanih
studenata po smeru i semestru, spisak nastavnika i dodele nastave, kao
i podrazumevana veličina grupe od 50 studenata (konstanta
\texttt{GROUP\_SIZE}), na osnovu koje se upis deli na približno
jednake grupe. Ravnomerna podela je pojednostavljenje: u stvarnosti
grupe nisu iste veličine, niti se studenti dele isključivo po smeru i
semestru. Ta konstanta se koristi samo u osnovnom modelu; u proširenom
modelu podela se preuzima iz dodela nastave, pa je tamo broj grupa
realniji, a veličine i dalje procenjene.

Te procene utiču na apsolutne vrednosti dveju veličina -- na broj
grupa, a time i na ukupan broj sesija, i na broj studenata koji u
mekom ograničenju kapaciteta (odeljak \ref{sec:kapacitet}) ostaju bez
mesta. Struktura problema time se ne menja, a poređenje dve
formulacije ostaje korektno, jer oba rešavača dobijaju isti ulaz.
Zaključci o tome kako se model ponaša pri rastu instance takođe ne
zavise od tačnosti procena, ali broj nenamirenih mesta u učionicama
treba čitati kao red veličine, a ne kao tačan podatak.

\subsection{Formiranje sesija}
\label{sec:sesije}

Pre nego što se pokrene rešavač, ulazni podaci se transformišu u
listu \emph{sesija} -- atomskih jedinica raspoređivanja. Ovaj
proces se odvija u modulu \texttt{data.py} i sastoji se od
sledećih koraka:

\begin{enumerate}
\item \textbf{Formiranje grupa studenata.}
  Grupa je kohorta koja zajedno sluša nastavu i formira se na jedan od
  tri načina, po prioritetu (funkcija \texttt{build\_groups} u modulu
  \code{data.py}):
  \begin{itemize}
  \item ako je zadat fajl sa osobljem u kome dodele nastave imenuju
    grupe (polje \texttt{groupIds}), imena grupa se uzimaju odatle, a
    njihove veličine iz broja upisanih -- tako radi realni ulaz i o tome
    detaljno govori odeljak \ref{sec:grupe_iz_osoblja};
  \item inače, ako ulazni JSON ima eksplicitnu \texttt{groups} sekciju,
    grupe se koriste onako kako su zadate;
  \item inače se za svaki par (smer, semestar) upis od $n$ studenata
    deli na $\lceil n / 50 \rceil$ grupa, gde je 50 podrazumevana
    veličina grupe (konstanta \texttt{GROUP\_SIZE}). Na primer, za
    Informatiku u 2. semestru sa 100 studenata formiraju se dve grupe.
  \end{itemize}

  Broj grupa zavisi, dakle, od izvora: nad punim letnjim semestrom
  deljenje po veličini daje 30 grupa, a imena iz dodela nastave 44
  grupe, jer stvarna podela na fakultetu ne prati pravilo od 50
  studenata. To se prenosi na broj sesija, pa instance osnovnog i
  proširenog modela nad istim ulaznim fajlom nisu iste veličine; o tome
  govore odeljci \ref{sec:grupe_iz_osoblja} i
  \ref{sec:finalna_metodologija}.

  Veličine grupa računaju se u oba slučaja iz broja upisanih, ali se
  ostatak pri deljenju tretira različito. Kod imenovanih grupa on se
  raspoređuje na prve grupe, pa je zbir veličina tačno jednak broju
  upisanih. Kod deljenja po veličini svaka grupa dobija $\lfloor n / k
  \rfloor$ studenata, gde je $k$ broj grupa, pa se ostatak gubi: za 55
  upisanih dobijaju se dve grupe po 27, a nad punim semestrom zbir
  veličina grupa iznosi 918 umesto 919. Na rezultate ove glave to ne
  utiče, jer osnovni model traži samo dopustivost i veličina grupe u
  njemu ne ulazi ni u jedno ograničenje; kapacitet učionice postaje
  meko ograničenje tek u proširenom modelu (odeljak
  \ref{sec:kapacitet}), gde je zbir veličina tačan.

\item \textbf{Dodeljivanje predmeta grupama.}
  Svaka grupa pohađa predmete koji odgovaraju njenom smeru i
  semestru.

\item \textbf{Generisanje sesija.}
  Za svaki predmet i svaku grupu, kreira se onoliko sesija koliko
  to nalaže kvota: \texttt{quota.theory} sesija tipa ``predavanje''
  i \texttt{quota.practice} sesija tipa ``vežbe''.
\end{enumerate}

Svaka sesija $s$ je okarakterisana sledećim atributima:
\begin{itemize}
\item identifikator grupe ($g_s$),
\item identifikator predmeta,
\item da li zahteva računare ($\text{comp}_s$),
\item tip časa (predavanje ili vežbe).
\end{itemize}

Za sesiju $s$, rešavač treba da odredi trojku $(d_s, h_s, r_s)$
-- dan, sat i učionicu u kojima će se čas održati.

\subsection{Obim modela i pretpostavke}
\label{sec:obim_modela}

Pošto se veći deo ulaza priprema pre pokretanja rešavača, korisno je
izričito reći šta rešavač odlučuje, a šta ne.

Rešavač odlučuje isključivo o trojki $(d_s, h_s, r_s)$ za svaku
sesiju. Van njegove nadležnosti ostaje sve ostalo: grupe studenata
formiraju se pre rešavanja, na jedan od tri opisana načina, i rešavač
ne raspoređuje pojedinačne studente u grupe niti menja njihov broj i
veličinu; nastavnici se preuzimaju iz unapred zadatih dodela nastave
(glava \ref{chp:osoblje}), pa rešavač ne bira ko drži koji čas, već
samo vodi računa da tako dodeljen nastavnik nema dva časa u istom
terminu; nedeljni broj časova po predmetu dolazi iz kvote, a sastav
kohorti koje zajedno slušaju predavanje iz dodela nastave.

Osnovna jedinica raspoređivanja je \emph{jednosatna sesija}. Predmet
sa tri časa vežbi nedeljno daje tri nezavisne sesije, a pravilo za
spajanje časova u blokove (odeljak \ref{sec:rule_join}) može zahtevati
da one budu u uzastopnim terminima istog dana, čime se dobija efekat
dvočasa ili tročasa. Ulazni model ima polje \texttt{duration}, ali ga
implementacija ne koristi i model pretpostavlja vrednost 1: vremenska
osa je diskretizovana na jednočasovne termine, a ograničenje
jedinstvenosti učionice poredi samo termine, ne i intervale. Podrška
za časove proizvoljnog trajanja zahtevala bi da promenljiva termina
nosi i trajanje i da se preklapanje sprečava ograničenjem nad
intervalima; to je naznačeno kao pravac daljeg rada.

\subsection{Nezavisnost predavanja i vežbi}
\label{sec:nezavisnost_tipova}

Naglašavamo da su \emph{predavanja i vežbe nezavisne jedinice
raspoređivanja}. Predmet u ulaznom modelu nije jedinica koja se
raspoređuje -- on samo povezuje časove koji pripadaju istoj nastavnoj
celini i propisuje njihov nedeljni broj kroz kvotu. Sve odluke o
raspoređivanju donose se na nivou pojedinačne sesije, pa se za
predavanja i vežbe istog predmeta nezavisno određuju:

\begin{itemize}
\item \textbf{broj časova} -- \texttt{quota.theory} i
  \texttt{quota.practice} su odvojena polja, tako da predmet može imati
  npr. 2 časa predavanja i 3 časova vežbi nedeljno ili ne mora uopšte imati vežbe.
\item \textbf{nastavnik} -- dodela nastave
  (\texttt{TeachingAssignment}, sekcija \ref{sec:staff_model}) vezuje se
  za par (predmet, tip nastave), pa predavanja i vežbe tipično drže
  različiti nastavnici,
\item \textbf{sastav slušalaca} -- kohorte se formiraju odvojeno po tipu
  (sekcija \ref{sec:cohorts}), što odgovara uobičajenoj praksi da više
  grupa sluša predavanje zajedno, a vežbe odvojeno,
\item \textbf{termin i učionica} -- promenljive
  $(d_s, h_s, r_s)$ postoje po sesiji, pa ni jedno ograničenje ne
  primorava predavanje i vežbe istog predmeta na isti termin ili istu
  učionicu,
\item \textbf{zahtev za računarima} -- $\text{comp}_s$ zavisi od tipa
  časa, o čemu govori nastavak.
\end{itemize}

Zahtev za računarima zaslužuje posebnu pažnju jer je u praksi najčešće
vezan isključivo za vežbe: predavanje iz programiranja drži se u
većim učionicama koje ne moraju imati računare, dok vežbe iz istog
predmeta zahtevaju računarsku učionicu.
Zbog toga se \texttt{needsComputers} ne zadaje kao jedna zastavica na
nivou predmeta, nego strukturom \texttt{ComputerNeed} sa odvojenim
poljima \texttt{theory} i \texttt{practice}. Predmet čije vežbe traže
računare, a predavanja ne, u JSON ulazu se zapisuje kao:

\begin{english}
\begin{lstlisting}[caption={Isečak ulaznog fajla: predmet sa računarima samo za vežbe}]
{"id": 4, "name": "Programiranje 1", "semester": 1, "trackId": 1,
 "quota": {"theory": 2, "practice": 3},
 "needsComputers": {"theory": false, "practice": true}}
\end{lstlisting}
\end{english}

Pri generisanju sesija svaka sesija preuzima vrednost koja odgovara
\emph{njenom} tipu nastave, a ne vrednost celog predmeta, pa je
$\text{comp}_s$ atribut sesije. Skalarni oblik
(\texttt{"needsComputers": true}) i dalje je dopušten i tumači se kao
isti zahtev za oba tipa nastave, radi kompatibilnosti sa starijim
ulaznim fajlovima.

\section{Formalna specifikacija}
\label{sec:formalna_specifikacija}

\subsection{Zajednička notacija}
\label{sec:notacija}

Obe formulacije opisuju isti problem nad istim ulazom, pa oznake
uvodimo na jednom mestu:

\begin{itemize}
\item $\mathcal{S}$ -- skup sesija, indeksiran sa
  $s \in \{0, 1, \ldots, |\mathcal{S}|-1\}$;
\item $D$ -- broj radnih dana, $H$ -- broj raspoloživih sati u toku
  jednog dana, $R$ -- broj učionica;
\item $\mathcal{S}_g \subseteq \mathcal{S}$ -- skup sesija koje pohađa
  grupa $g$. Sesija koju zajednički sluša više grupa pripada skupu
  svake od njih (odeljak \ref{sec:cohorts});
\item $\text{comp}_s$ -- zastavica koja kaže da li sesija $s$ zahteva
  računarsku učionicu. To je atribut sesije, a ne predmeta (odeljak
  \ref{sec:nezavisnost_tipova});
\item $\text{hasComputers}(r)$ -- da li učionica $r$ ima računare;
\item $\text{Elig}(s) = \{r : \neg\,\text{comp}_s \lor
  \text{hasComputers}(r)\}$ -- skup učionica u kojima sesija $s$ sme
  da se održi.
\end{itemize}

Rešenje pridružuje svakoj sesiji $s$ trojku $(d_s, h_s, r_s)$ -- dan,
sat i učionicu. Zahtev za računarima jedini je tvrdi uslov na izbor
učionice: kapacitet ne ulazi u $\text{Elig}(s)$, već se tek u glavi
\ref{chp:nadogradnja} uvodi kao meko ograničenje (odeljak
\ref{sec:kapacitet}). Osnovni model radi u režimu izvodljivosti --
traži se bilo koje rešenje koje zadovoljava sva tvrda ograničenja, bez
funkcije cilja.

\subsection{CP model: promenljive, domeni i ograničenja}
\label{sec:cp_form}

\textbf{Promenljive i domeni.}
Za svaku sesiju $s \in \{0, 1, \ldots, |\mathcal{S}|-1\}$ definišemo
sledeće celobrojne promenljive:

\begin{align}
\text{day}_s &\in \{0, 1, \ldots, D-1\} \label{eq:cp_day} \\
\text{slot}_s &\in \{0, 1, \ldots, H-1\} \label{eq:cp_slot} \\
\text{room}_s &\in \{0, 1, \ldots, R-1\} \label{eq:cp_room}
\end{align}

Dodatno, definišemo dve izvedene promenljive:

\begin{align}
\text{flatTime}_s &= \text{day}_s \cdot H + \text{slot}_s
  \label{eq:cp_flat} \\
\text{roomTime}_s &= \text{room}_s \cdot (D \cdot H) +
  \text{flatTime}_s \label{eq:cp_roomtime}
\end{align}

Promenljiva $\text{flatTime}_s$ linearizuje par (dan, sat) u
jedinstven indeks unutar nedelje, a $\text{roomTime}_s$ trojku
(učionica, dan, sat) u jedinstven indeks preko svih učionica i
termina. Izvedene promenljive nisu automatski vezane za osnovne: one
se uvode kao nove celobrojne promenljive, a jednakosti
\eqref{eq:cp_flat} i \eqref{eq:cp_roomtime} dodaju se eksplicitno kao
ograničenja. Ukupno model ima $5 \cdot |\mathcal{S}|$ celobrojnih
promenljivih.

\medskip
\textbf{HC1: Jedinstvenost učionice u terminu.}
Nikoje dve sesije ne mogu se održavati u istoj učionici u isto vreme:

\begin{equation}
\texttt{AllDifferent}(\text{roomTime}_0, \text{roomTime}_1,
  \ldots, \text{roomTime}_{|\mathcal{S}|-1})
\label{eq:cp_hc1}
\end{equation}

Ovo globalno ograničenje obuhvata sve sesije i garantuje
da su sve trojke $(d_s, h_s, r_s)$ međusobno različite.

\medskip
\textbf{HC2: Bez kolizija unutar grupe.}
Nikoje dve sesije iste grupe ne smeju biti u isto vreme:

\begin{equation}
\texttt{AllDifferent}(\{\text{flatTime}_s : s \in \mathcal{S}_g\}),
  \quad \forall g
\label{eq:cp_hc2}
\end{equation}

Oba ograničenja imaju isti oblik, samo nad različitim skupovima
promenljivih: HC1 nad svim vrednostima $\text{roomTime}$, a HC2 nad
vrednostima $\text{flatTime}$ jedne grupe. Sesija ulazi u HC2 za
\emph{svaku} grupu kojoj pripada, što je bitno za zajedničke sesije
koje istovremeno pohađa više grupa.

\medskip
\textbf{HC3: Ograničenje računarskih učionica.}
Ako sesija $s$ zahteva računare, njen izbor učionice je ograničen
na podskup učionica sa računarima:

\begin{equation}
\text{comp}_s \implies \text{room}_s \in
  \texttt{AllowedAssignments}(\{r : \text{hasComputers}(r)\})
\label{eq:cp_hc3}
\end{equation}

Pošto je $\text{comp}_s$ atribut sesije, ograničenje se odnosi samo na
one sesije čiji tip nastave traži računare: predavanja iz predmeta
čije vežbe traže računare ostaju slobodna da idu u bilo koju učionicu.

\subsection{MIP model: promenljive, domeni i ograničenja}
\label{sec:mip_form}

\textbf{Promenljive i domeni.}
Za svaku sesiju $s$ i svaku dopustivu kombinaciju (dan, sat, učionica)
definišemo binarnu promenljivu:

\begin{equation}
x_{s,d,h,r} \in \{0, 1\}, \quad \forall s, \forall d \in
  \{0,\ldots,D\!-\!1\}, \forall h \in \{0,\ldots,H\!-\!1\},
  \forall r \in \text{Elig}(s)
\label{eq:mip_vars}
\end{equation}

\noindent Promenljive za učionice van skupa $\text{Elig}(s)$ se uopšte
ne kreiraju, čime se ograničenje računarskih učionica poštuje
implicitno, kroz domen modela, a ne kroz zasebno ograničenje. Ukupan
broj promenljivih je $O(|\mathcal{S}| \cdot D \cdot H \cdot R)$.

Sve promenljive modela su binarne, a sva ograničenja linearna, pa je
ovo binarni celobrojni linearni model (0-1 ILP), poseban slučaj opšteg
MIP oblika iz odeljka \ref{sec:mip}. Rešava ga SCIP, odakle i oznaka
MIP/SCIP koju koristimo u tabelama.

\medskip
\textbf{HC1: Svaka sesija raspoređena tačno jednom.}

\begin{equation}
\sum_{d,h,r} x_{s,d,h,r} = 1, \quad \forall s \in \mathcal{S}
\label{eq:mip_hc1}
\end{equation}

\medskip
\textbf{HC2: Najviše jedna sesija u datoj učionici u datom terminu.}

\begin{equation}
\sum_{s \in \mathcal{S}} x_{s,d,h,r} \leq 1, \quad \forall d,h,r
\label{eq:mip_hc2}
\end{equation}

\medskip
\textbf{HC3: Bez kolizija unutar grupe.}
Za svaku grupu $g$ i svaki termin $(d,h)$:

\begin{equation}
\sum_{s \in \mathcal{S}_g} \sum_{r} x_{s,d,h,r} \leq 1,
  \quad \forall g, \forall d, \forall h
\label{eq:mip_hc3}
\end{equation}

Kao i CP model iz prethodnog pododeljka, osnovni MIP model radi u
režimu izvodljivosti, bez funkcije cilja. Meka ograničenja i funkcija
cilja uvode se tek u proširenom CP modelu (glava
\ref{chp:nadogradnja}).

\subsection{Ekvivalentnost dve formulacije}
\label{sec:ekvivalentnost}

Ograničenja dva modela odgovaraju jedno drugom. MIP ograničenje HC1
obezbeđuje da svaka sesija dobije tačno jednu dodelu; u CP modelu to
važi po konstrukciji, jer svaka sesija ima po jednu promenljivu za dan,
sat i učionicu. MIP ograničenje HC2 odgovara ograničenju
\texttt{AllDifferent} nad $\text{roomTime}$, a HC3 ograničenju
\texttt{AllDifferent} nad $\text{flatTime}$ po grupi. Zahtev za
računarima oba modela izražavaju kroz isti skup $\text{Elig}(s)$ --
CP eksplicitnim ograničenjem nad domenom promenljive
$\text{room}_s$, a MIP izostavljanjem promenljivih za nedozvoljene
učionice.

Obe formulacije zato opisuju isti skup dopustivih rešenja. To je uslov
bez koga poređenje iz glave \ref{chp:evaluacija} ne bi bilo
smisleno: razlike u izmerenim veličinama potiču od načina zapisa
problema i od rešavača, a ne od toga da jedan model rešava lakši
problem.

\section{Implementacija}
\label{sec:implementacija}

Obe specifikacije prevedene su u kod korišćenjem biblioteke OR-Tools:
CP model u klasi \texttt{SimpleCPSolver} (modul \code{cp_solver.py}),
a MIP model u klasi \texttt{SimpleMIPSolver} (modul
\code{mip_solver.py}). U nastavku prikazujemo isečke koji pokazuju
kako se ulazni podaci i elementi specifikacije iz prethodnog odeljka
zapisuju u kodu.

\subsection{Ulazni model u kodu}

Entiteti opisani u odeljku \ref{sec:input_model} definisani su kao
Pydantic klase, sa tipovima koji služe i kao specifikacija ulaza i kao
provera koju biblioteka sprovodi pri učitavanju:

\begin{english}
\begin{lstlisting}[language=Python, caption={Model ulaznih podataka}, label={lst:pydantic_model}]
@dataclass(config=_config)
class Settings:
    working_days: List[str] = Field(alias="workingDays")
    start_hour: int = Field(alias="startHour")
    end_hour: int = Field(alias="endHour")
    duration: int = 1

@dataclass(config=_config)
class Classroom:
    id: int
    name: str
    loc_id: int = Field(alias="locId")
    has_computers: bool = Field(alias="hasComputers")
    capacity: int = 0

@dataclass(config=_config)
class SchedulingInput:
    settings: Settings
    locations: List[Location]
    classrooms: List[Classroom]
    courses: List[Course]
    tracks: List[StudyTrack] = Field(
        validation_alias=AliasChoices("tracks", "departments"))
    students_enrolled: List[StudentsEnrolled] = Field(
        alias="studentsEnrolled")
    rules: dict[str, RuleConfig] = Field(default_factory=dict)
    groups: List[GroupDef] = Field(default_factory=list)
\end{lstlisting}
\end{english}

\subsection{CP implementacija}

Izvedene promenljive $\text{flatTime}_s$ i $\text{roomTime}_s$ uvode se
kao nove celobrojne promenljive, a jednakosti \eqref{eq:cp_flat} i
\eqref{eq:cp_roomtime} dodaju se eksplicitno:

\begin{english}
\begin{lstlisting}[language=Python, caption={Linearizacija termina i para (učionica, termin)}]
self.flat_time_var[s] = self.model.NewIntVar(
    0, total_slots - 1, f"flat_time_{s}")
self.model.Add(
    self.flat_time_var[s]
    == self.day_var[s] * H + self.slot_var[s])

self.room_time_var[s] = self.model.NewIntVar(
    0, R * total_slots - 1, f"room_time_{s}")
self.model.Add(
    self.room_time_var[s]
    == self.room_var[s] * total_slots + self.flat_time_var[s])
\end{lstlisting}
\end{english}

Ograničenja HC1 i HC2 svode se na isti poziv nad različitim skupovima
promenljivih:

\begin{english}
\begin{lstlisting}[language=Python, caption={Ograničenja HC1 i HC2 u CP-SAT modelu}, label={lst:cp_alldiff}]
# HC1: nikoje dve sesije ne dele istu (ucionicu, dan, sat)
self.model.AddAllDifferent(list(self.room_time_var.values()))

# HC2: nikoje dve sesije iste grupe nisu u istom terminu
groups = defaultdict(list)
for s, session in enumerate(self.sessions):
    for group_id in session.group_ids:
        groups[group_id].append(s)

for group_id, session_indices in groups.items():
    self.model.AddAllDifferent(
        [self.flat_time_var[s] for s in session_indices])
\end{lstlisting}
\end{english}

Uz pet promenljivih iz specifikacije, implementacija po sesiji gradi i
šestu, $\text{loc}_s$, koja se iz $\text{room}_s$ dobija indeksiranjem
niza lokacija. Ona ne utiče na skup dopustivih rešenja osnovnog
modela, ali je koriste pravila vezana za lokacije u glavama
\ref{chp:nadogradnja} i \ref{chp:osoblje}, pa je izmereni broj
promenljivih u glavi \ref{chp:evaluacija} $6 \cdot |\mathcal{S}|$.

Skup $\text{Elig}(s)$ iz ograničenja HC3 gradi se po sesiji, a
ograničenje se dodaje samo kada je taj skup uži od skupa svih
učionica:

\begin{english}
\begin{lstlisting}[language=Python, caption={Dozvoljene učionice po sesiji}]
for s, session in enumerate(self.sessions):
    allowed = [
        i for i, room in enumerate(self.classrooms)
        if not session.needs_computers or room.has_computers
    ]
    if len(allowed) < len(self.classrooms):
        self.model.AddAllowedAssignments(
            [self.room_var[s]], [[idx] for idx in allowed])
\end{lstlisting}
\end{english}

\subsection{MIP implementacija}

Binarne promenljive $x_{s,d,h,r}$ kreiraju se samo za učionice iz
skupa $\text{Elig}(s)$, koji vraća pomoćna metoda
\texttt{\_eligible\_room\_indices}:

\begin{english}
\begin{lstlisting}[language=Python, caption={Kreiranje binarnih promenljivih}]
for s, session in enumerate(self.sessions):
    self.x[s] = {}
    eligible = self._eligible_room_indices(session)
    for d in range(D):
        for h in range(H):
            for r in eligible:
                self.x[s][(d, h, r)] = self.solver.BoolVar(
                    f"x_{s}_{d}_{h}_{r}")
\end{lstlisting}
\end{english}

Svako od tri ograničenja zapisuje se kao familija linearnih
nejednakosti, po jedna za svaku kombinaciju indeksa nad kojima se
kvantifikuje; rečnik \texttt{groups} gradi se isto kao u CP modelu:

\begin{english}
\begin{lstlisting}[language=Python, caption={Tvrda ograničenja u MIP modelu}]
# HC1: svaka sesija rasporedjena tacno jednom
for s in range(S):
    self.solver.Add(sum(self.x[s].values()) == 1)

# HC2: najvise jedna sesija u datoj (dan, sat, ucionica)
for d in range(D):
    for h in range(H):
        for r in range(R):
            vars_at_dhr = [
                self.x[s][(d, h, r)]
                for s in range(S) if (d, h, r) in self.x[s]
            ]
            if len(vars_at_dhr) > 1:
                self.solver.Add(sum(vars_at_dhr) <= 1)

# HC3: jedna grupa nema dve sesije u istom terminu
for group_id, session_indices in groups.items():
    for d in range(D):
        for h in range(H):
            group_vars = [
                self.x[s][(d, h, r)]
                for s in session_indices for r in range(R)
                if (d, h, r) in self.x[s]
            ]
            if len(group_vars) > 1:
                self.solver.Add(sum(group_vars) <= 1)
\end{lstlisting}
\end{english}

Razlika u načinu zapisa vidi se već ovde: ono što CP model izražava
jednim pozivom nad skupom promenljivih, MIP model zapisuje petljom
koja generiše po jednu nejednakost za svaku kombinaciju indeksa.

\section{Poređenje formulacija}
\label{sec:poredjenje_formulacija}

Tabela \ref{tbl:formulacije} sumira ključne razlike između dve
formulacije.

\begin{table}[ht]
\centering
\caption{Poređenje CP i MIP formulacija}
\label{tbl:formulacije}
\begin{tabular}{lll}
\toprule
\textbf{Aspekt} & \textbf{CP-SAT} & \textbf{MIP/SCIP} \\
\midrule
Tip promenljivih & Celobrojne & Binarne \\
Broj promenljivih & $O(|\mathcal{S}|)$ & $O(|\mathcal{S}| \cdot D \cdot H \cdot R)$ \\
Tip ograničenja & Globalna & Linearne nejednakosti \\
Broj ograničenja & Kompaktan & $O(D \cdot H \cdot R + G \cdot D \cdot H)$ \\
Pretraga & Propagacija + SAT & LP relaksacija + B\&B \\
\bottomrule
\end{tabular}
\end{table}

CP formulacija koristi mali broj celobrojnih promenljivih sa globalnim ograničenjima,
dok MIP formulacija eksplicitno enumeriše sve moguće dodele kroz
binarne promenljive. Ovo ima uticaj na memorijski otisak
i brzinu konstrukcije modela.

Kako se te razlike ponašaju na realnim podacima ne može se zaključiti
iz same formulacije, pa u narednoj glavi obe implementacije merimo na
tri skale problema.


% ------------------------------------------------------------------------------
\chapter{Preliminarna evaluacija osnovnog modela}
\label{chp:evaluacija}
\chaptermark{Preliminarna evaluacija}
% ------------------------------------------------------------------------------

U ovoj glavi opisujemo metodologiju merenja performansi, predstavljamo
rezultate eksperimenata i analiziramo dobijene podatke.

Evaluacija je preliminarna u dva smisla. Odnosi se na osnovni model iz
glave \ref{chp:formulacija}, dakle na traženje dopustivog rasporeda uz
tri tvrda ograničenja i bez funkcije cilja, i služi tome da se izabere
formulacija koju ćemo dalje razvijati. Završna evaluacija proširenog
modela, sa mekim ograničenjima, funkcijom cilja i nastavnim osobljem,
data je odvojeno u glavi \ref{chp:finalna-evaluacija}.

\section{Metodologija}
\label{sec:metodologija}

\subsection{Podskupovi realnih podataka}

Eksperimenti su sprovedeni na tri podskupa realnih podataka letnjeg
semestra (\code{input_full_2_semester.json}), dobijena filtriranjem po
godinama studija:

\begin{enumerate}
\item \textbf{MATF-S} (mali) -- 1. godina (semestar 2): 30 predmeta,
  8 grupa i 204 sesije.

\item \textbf{MATF-M} (srednji) -- 1. i 2. godina (semestri 2 i 4):
  63 predmeta, 16 grupa i 398 sesija.

\item \textbf{MATF-L} (veliki) -- sve četiri godine (semestri 2, 4, 6
  i 8): 134 predmeta, 30 grupa i 722 sesije. Ovo je pun obim nastave
  jednog semestra.
\end{enumerate}

Skalira se, dakle, samo potražnja: sve tri instance zadržavaju svih 29
učionica na tri lokacije, pa se razlikuju jedino po količini nastave
koju treba smestiti. Podskupovi su namerno definisani isto kao u
završnoj evaluaciji (odeljak \ref{sec:finalna_metodologija}) i nad
istim ulaznim fajlom, kako bi se rezultati pre i posle proširenja
modela odnosili na istu instancu.

Jedna razlika ipak ostaje i treba je imati u vidu. Ovde se grupe
studenata dobijaju deljenjem broja upisanih na grupe od najviše 50
studenata, jer osnovni model ne poznaje nastavno osoblje, dok se u
glavi \ref{chp:finalna-evaluacija} grupe izvode iz imena navedenih u
dodelama nastave (odeljak \ref{sec:grupe_iz_osoblja}), koja opisuju
stvarnu podelu na fakultetu. Na najvećoj skali to je razlika između 30
i 44 grupe, pa ista instanca tamo ima 839 umesto 722 sesije. Sam
postupak raspoređivanja se ne menja -- menja se samo koliko časova
ulazi u model -- ali apsolutne vrednosti iz ove glave i iz glave
\ref{chp:finalna-evaluacija} zato nisu neposredno uporedive.

\subsection{Metrike}

Za svaki par (rešavač, skala) merimo sledeće metrike:

\begin{itemize}
\item \textbf{Broj sesija} -- ukupan broj časova koje treba rasporediti.
\item \textbf{Broj promenljivih} -- veličina modela u terminima
  odlučivačkih promenljivih.
\item \textbf{Broj ograničenja} -- ukupan broj ograničenja u modelu.
\item \textbf{Vreme konstrukcije} -- vreme potrebno za kreiranje
  modela (promenljivih i ograničenja), mereno pomoću
  \texttt{time.perf\_counter()}.
\item \textbf{Vreme rešavanja} -- vreme potrebno za pronalaženje
  dopustivog rešenja.
\item \textbf{Ukupno vreme} -- zbir vremena konstrukcije i rešavanja.
\item \textbf{Memorija modela} -- memorijski otisak modela, meren
  pomoću \texttt{tracemalloc} (razlika memorijskih snimaka pre i
  posle konstrukcije modela).
\item \textbf{Maksimalna RAM} -- vršna potrošnja memorije procesa,
  merena pomoću \texttt{resource.getrusage()}.
\item \textbf{Status} -- da li je rešavač našao dopustivo rešenje
  ili je isteklo vreme.
\item \textbf{Validnost rešenja} -- nezavisna provera da li dobijeno
  rešenje zadovoljava sva tvrda ograničenja (bez kolizija učionica,
  bez kolizija grupa, poštovanje računarskih zahteva).
\end{itemize}

\subsection{Validacija rešenja}

Rezultati svakog rešavača se nezavisno verifikuju funkcijom
\texttt{validate\_solution} koja proverava:
\begin{enumerate}
\item da nijedna dva časa ne dele istu trojku (dan, sat, učionica),
\item da nijedna dva časa iste grupe nisu u istom terminu (dan, sat),
\item da časovi koji zahtevaju računare budu u učionicama sa
  računarima (provera se vrši nad zahtevom sesije, koji zavisi od tipa
  nastave),
\item da broj dodeljenih termina odgovara broju sesija.
\end{enumerate}

\section{Rezultati}
\label{sec:rezultati}

U tabelama \ref{tbl:results_s}, \ref{tbl:results_m} i
\ref{tbl:results_l} prikazani su rezultati za svaku skalu.

\begin{table}[ht]
\centering
\caption{Rezultati za MATF-S (1. godina)}
\label{tbl:results_s}
\begin{tabular}{lrr}
\toprule
\textbf{Metrika} & \textbf{CP-SAT} & \textbf{MIP/SCIP} \\
\midrule
Sesije                 & \multicolumn{2}{c}{204} \\
Promenljive            & $1\,224$        & $319\,680$ \\
Ograničenja            & $649$           & $2\,424$ \\
Vreme konstrukcije     & 0,016 s         & 5,121 s \\
Vreme rešavanja        & 0,026 s         & 3,206 s \\
Ukupno vreme           & 0,042 s         & 8,327 s \\
Memorija modela        & 0,22 MB         & 71,54 MB \\
Status                 & \textsc{feasible} & \textsc{feasible} \\
Rešenje validno        & da              & da \\
\bottomrule
\end{tabular}
\end{table}

\begin{table}[ht]
\centering
\caption{Rezultati za MATF-M (1--2. godina)}
\label{tbl:results_m}
\begin{tabular}{lrr}
\toprule
\textbf{Metrika} & \textbf{CP-SAT} & \textbf{MIP/SCIP} \\
\midrule
Sesije                 & \multicolumn{2}{c}{398} \\
Promenljive            & $2\,388$        & $619\,440$ \\
Ograničenja            & $1\,269$        & $3\,098$ \\
Vreme konstrukcije     & 0,030 s         & 9,879 s \\
Vreme rešavanja        & 0,084 s         & 8,170 s \\
Ukupno vreme           & 0,114 s         & 18,049 s \\
Memorija modela        & 0,39 MB         & 138,56 MB \\
Status                 & \textsc{feasible} & \textsc{feasible} \\
Rešenje validno        & da              & da \\
\bottomrule
\end{tabular}
\end{table}

\begin{table}[ht]
\centering
\caption{Rezultati za MATF-L (sve četiri godine)}
\label{tbl:results_l}
\begin{tabular}{lrr}
\toprule
\textbf{Metrika} & \textbf{CP-SAT} & \textbf{MIP/SCIP} \\
\midrule
Sesije                 & \multicolumn{2}{c}{722} \\
Promenljive            & $4\,332$        & $1\,102\,560$ \\
Ograničenja            & $2\,319$        & $4\,262$ \\
Vreme konstrukcije     & 0,055 s         & 18,600 s \\
Vreme rešavanja        & 1,598 s         & 17,804 s \\
Ukupno vreme           & 1,652 s         & 36,404 s \\
Memorija modela        & 0,70 MB         & 246,57 MB \\
Status                 & \textsc{feasible} & \textsc{feasible} \\
Rešenje validno        & da              & da \\
\bottomrule
\end{tabular}
\end{table}

\textbf{Napomena:} Prikazani rezultati su dobijeni pokretanjem
benchmark-a komandom:

\begin{english}
\begin{verbatim}
bazel run //src/algo:benchmark -- --mode compare --json report.json
\end{verbatim}
\end{english}

\section{Analiza rezultata}
\label{sec:analiza}

Na osnovu sprovedenih eksperimenata, možemo izvući sledeće zaključke:

\subsection{Kompaktnost modela}

CP-SAT formulacija koristi $O(|\mathcal{S}|)$ celobrojnih
promenljivih, dok MIP/SCIP formulacija koristi
$O(|\mathcal{S}| \cdot D \cdot H \cdot R)$ binarnih promenljivih.
Sa $D=5$ dana, $H=12$ sati i $R=29$ učionica svaka sesija u MIP-u
generiše do $5 \cdot 12 \cdot 29 = 1740$ binarnih promenljivih,
naspram šest celobrojnih u CP-u. Izmereni broj od šest promenljivih po
sesiji je za jednu veći od pet iz specifikacije (odeljak
\ref{sec:cp_form}), jer implementacija uz njih uvek gradi i pomoćnu
promenljivu lokacije, koju koriste pravila iz glava
\ref{chp:nadogradnja} i \ref{chp:osoblje}. Ova razlika od dva do tri
reda veličine odražava se na memorijski otisak i vreme konstrukcije
modela.

\subsection{Brzina rešavanja}

CP-SAT rešavač koristi globalna ograničenja
(\texttt{AllDifferent}, \texttt{AllowedAssignments}) sa
specijalizovanim algoritmima propagacije koji rano eliminišu
nedopustive vrednosti iz domena promenljivih. U kombinaciji sa
SAT tehnologijama učenja konflikata, ovo omogućava znatno bržu
pretragu prostora rešenja u poređenju sa LP relaksacijom i
branch-and-bound pristupom MIP/SCIP rešavača.

\subsection{Skalabilnost}

Razlika u performansama između dva pristupa postaje izraženija sa
rastom veličine problema. Oba rešavača su na sve tri skale našla
validan raspored u okviru zadatog limita, ali MIP/SCIP-u je za to
potrebno znatno više vremena, i to vreme raste brže: od 8,3 sekunde na
MATF-S do 36,4 sekunde na MATF-L, naspram 0,04 i 1,65 sekundi kod
CP-SAT-a. Pošto sve tri instance koriste isti skup od 29 učionica,
razlika potiče isključivo od porasta broja sesija.

\subsection{Teorijska kompleksnost}

Tabela \ref{tbl:complexity} sumira teorijsku kompleksnost
oba pristupa.

\begin{table}[ht]
\centering
\caption{Teorijska kompleksnost formulacija}
\label{tbl:complexity}
\begin{tabular}{lll}
\toprule
\textbf{Aspekt} & \textbf{CP-SAT} & \textbf{MIP/SCIP} \\
\midrule
Promenljive & $O(|\mathcal{S}|)$ &
  $O(|\mathcal{S}| \cdot D \cdot H \cdot R)$ \\
Ograničenja & Kompaktna globalna &
  $O(D \cdot H \cdot R + G \cdot D \cdot H)$ \\
Propagacija & Specijalizovani algoritmi &
  LP relaksacija \\
Učenje & Konflikti (SAT klauze) &
  Ravni odsecanja \\
\bottomrule
\end{tabular}
\end{table}

Rezultati ove glave odnose se na osnovni model i služe izboru
formulacije za dalji razvoj. Pošto se model u glavama
\ref{chp:nadogradnja} i \ref{chp:osoblje} znatno proširuje, u glavi
\ref{chp:finalna-evaluacija} dajemo završnu evaluaciju konačne verzije
rešavača na istim trima skalama.

U narednoj glavi zato prelazimo sa poređenja dve formulacije na razvoj
jedne: CP model dobija arhitekturu pravila, funkciju cilja i prva
ograničenja kvaliteta rasporeda.


% ------------------------------------------------------------------------------
\chapter{Nadogradnja CP rešavača dodatnim ograničenjima}
\label{chp:nadogradnja}
\chaptermark{Nadogradnja CP rešavača}
% ------------------------------------------------------------------------------

Na osnovu rezultata iz glave \ref{chp:evaluacija}, CP-SAT rešavač je
izabran kao osnova za dalji razvoj sistema. U ovoj glavi opisujemo
nadogradnju modela dodatnim ograničenjima koja raspored približavaju
realnim potrebama Matematičkog fakulteta. Za razliku od prvog dela rada,
gde smo tražili bilo koje dopustivo rešenje, ovde uvodimo i funkciju
cilja koja optimizuje kvalitet rasporeda kroz meka ograničenja.

Implementirana su četiri nova ograničenja:
\begin{itemize}
\item \textbf{spajanje časova u blokove} -- časovi istog predmeta, grupe
  i tipa grupišu se u uzastopne blokove (na primer dvočasi i tročasi);
\item \textbf{jedna lokacija po danu za grupu} -- grupa studenata ne menja
  zgradu u toku jednog radnog dana;
\item \textbf{bez procepa u rasporedu} -- izbegavanje slobodnih sati
  između dva zauzeta sata u istom danu;
\item \textbf{kapacitet učionice} -- prekoračenje kapaciteta se kažnjava
  brojem studenata bez mesta, ali nije zabranjeno.
\end{itemize}

\section{Arhitektura pravila}
\label{sec:arhitektura_pravila}

Svako ograničenje je modelovano kao zaseban dodatak (eng. \emph{plugin})
koji se priključuje rešavaču. Ovakva arhitektura omogućava da se pravila
nezavisno uključuju, isključuju i konfigurišu, bez izmena u jezgru
rešavača.

Osnovu čini apstraktna klasa \texttt{SchedulingRule} (modul
\texttt{src/rules/base.py}):

\begin{english}
\begin{lstlisting}[language=Python, caption={Apstraktna klasa pravila}]
class SchedulingRule(ABC):
    def __init__(self, enabled=True, penalty=0, **kwargs):
        self.enabled = enabled
        self.penalty = penalty

    @property
    def is_hard(self) -> bool:
        return self.penalty == 0

    @abstractmethod
    def apply(self, solver) -> list[cp_model.IntVar]:
        ...
\end{lstlisting}
\end{english}

Svako pravilo ima dva parametra: \texttt{enabled} (da li je pravilo
aktivno) i \texttt{penalty} (cena kršenja). Vrednost \texttt{penalty}
istovremeno određuje i \emph{tip} ograničenja:

\begin{itemize}
\item \texttt{penalty == 0} -- \textbf{tvrdo} (eng. \emph{hard})
  ograničenje. Mora biti zadovoljeno; u suprotnom rešavač prijavljuje
  da je problem nedopustiv (\texttt{INFEASIBLE}). Metoda \texttt{apply}
  dodaje ograničenja direktno u model i vraća praznu listu.
\item \texttt{penalty > 0} -- \textbf{meko} (eng. \emph{soft})
  ograničenje. Rešavač pokušava da ga zadovolji, ali ga može prekršiti
  uz datu cenu. Metoda \texttt{apply} vraća listu binarnih promenljivih
  prekršaja (vrednost 1 označava prekršaj).
\end{itemize}

\subsection{Konfiguracija pravila}

Pravila se konfigurišu u ulaznom JSON fajlu pomoću sekcije
\texttt{rules}, gde svaki ključ odgovara nazivu pravila:

\begin{english}
\begin{lstlisting}[language=Python, caption={Primer konfiguracije pravila u JSON-u}]
"rules": {
    "joinSameClasses":  { "enabled": true, "penalty": 0 },
    "noGapsInSchedule": { "enabled": true, "penalty": 5 }
}
\end{lstlisting}
\end{english}

Ova konfiguracija se učitava u Pydantic model \texttt{RuleConfig}
(modul \texttt{model.py}). U gornjem primeru, pravilo
\texttt{joinSameClasses} je tvrdo (penalty 0), dok je
\texttt{noGapsInSchedule} meko sa cenom 5 po prekršaju.

\subsection{Primena pravila i funkcija cilja}

Rešavač održava registar pravila \texttt{RULE\_REGISTRY} koji preslikava
naziv pravila u odgovarajuću klasu. Metoda \texttt{apply\_rules}
instancira sva uključena pravila, prikuplja promenljive prekršaja i
formira funkciju cilja:

\begin{english}
\begin{lstlisting}[language=Python, caption={Primena pravila i formiranje funkcije cilja}]
def apply_rules(self):
    self.penalty_vars = []
    for rule_name, config in self.scheduling_input.rules.items():
        if not config.enabled:
            continue
        rule_class = RULE_REGISTRY.get(rule_name)
        rule = rule_class(enabled=config.enabled, penalty=config.penalty)
        violations = rule.apply(self)
        for v in violations:
            self.penalty_vars.append((config.penalty, v))

    if self.penalty_vars:
        self.model.Minimize(
            sum(weight * var for weight, var in self.penalty_vars)
        )
\end{lstlisting}
\end{english}

Funkcija cilja je dakle težinska suma svih prekršaja mekih ograničenja:

\begin{equation}
\min \sum_{r} \text{penalty}_r \cdot \sum_{v \in V_r} v
\label{eq:objective}
\end{equation}

gde je $V_r$ skup promenljivih prekršaja koje vraća pravilo $r$, a
$\text{penalty}_r$ njegova cena. Tvrda ograničenja ne učestvuju u
funkciji cilja jer vraćaju praznu listu prekršaja.

\section{CP-SAT tehnike modeliranja}
\label{sec:cpsat_tehnike}

Nova pravila koriste nekoliko naprednih tehnika modeliranja koje pruža
CP-SAT \cite{cpsat2020}. U ovom odeljku ih ukratko objašnjavamo, jer se
koriste u sve tri implementacije.

\subsection{Reifikacija i \texttt{OnlyEnforceIf}}

\emph{Reifikacija} povezuje istinitost nekog ograničenja sa binarnom
\emph{indikatorskom} promenljivom (literalom). U CP-SAT, metoda
\texttt{OnlyEnforceIf(b)} čini da ograničenje važi \emph{samo ako} je
literal $b$ tačan:

\begin{english}
\begin{lstlisting}[language=Python, caption={Uslovno (reifikovano) ograničenje}]
# lit == 1  =>  flat_time[s] == t
model.Add(flat_time[s] == t).OnlyEnforceIf(lit)
# lit == 0  =>  flat_time[s] != t
model.Add(flat_time[s] != t).OnlyEnforceIf(lit.Not())
\end{lstlisting}
\end{english}

Navođenjem oba smera (za $b$ i za $\neg b$, tj. \texttt{lit.Not()})
dobija se puna ekvivalencija $b \Leftrightarrow (\text{flat\_time}_s = t)$,
čime literal $b$ tačno odražava istinitost ograničenja.

\subsection{Logičko ILI preko \texttt{AddMaxEquality}}

Za binarne promenljive (vrednosti 0 ili 1), maksimum je jednak logičkom
ILI. CP-SAT to izražava ograničenjem \texttt{AddMaxEquality}:

\begin{english}
\begin{lstlisting}[language=Python, caption={Logičko ILI nad binarnim promenljivama}]
# b == 1 ako je bar jedan od literala u "lits" jednak 1
model.AddMaxEquality(b, lits)   # b = max(lits) = OR(lits)
\end{lstlisting}
\end{english}

Ovo koristimo da utvrdimo da li je grupa zauzeta u nekom terminu (bar
jedna sesija je u tom terminu).

\subsection{\texttt{AddBoolAnd} i \texttt{AddBoolOr}}

Ograničenja \texttt{AddBoolAnd} (konjunkcija) i \texttt{AddBoolOr}
(disjunkcija) nameću da svi, odnosno bar jedan literal budu tačni. U
kombinaciji sa \texttt{OnlyEnforceIf} koriste se za potpunu reifikaciju
složenih logičkih izraza:

\begin{english}
\begin{lstlisting}[language=Python, caption={Reifikacija ekvivalencije logičkog izraza}]
# gap == 1  <=>  (has_before AND has_after AND NOT busy)
model.AddBoolAnd([has_before, has_after, busy.Not()]).OnlyEnforceIf(gap)
model.AddBoolOr([has_before.Not(), has_after.Not(), busy]).OnlyEnforceIf(gap.Not())
\end{lstlisting}
\end{english}

Prvi red obezbeđuje smer $gap \Rightarrow (\ldots)$, a drugi
(primenom De Morganovih zakona) smer $\neg gap \Rightarrow \neg(\ldots)$.

\subsection{Indeksiranje niza preko \texttt{AddElement}}

Ograničenje \texttt{AddElement(index, array, target)} nameće
$\text{target} = \text{array}[\text{index}]$, gde je \texttt{index}
promenljiva. Ovo omogućava ``pretragu'' kroz niz na osnovu vrednosti
druge promenljive:

\begin{english}
\begin{lstlisting}[language=Python, caption={Indeksiranje niza promenljivom}]
# loc[s] == locations_of_classrooms[room_var[s]]
model.AddElement(room_var[s], locations_of_classrooms, loc[s])
\end{lstlisting}
\end{english}

\section{Pravilo: spajanje časova u blokove}
\label{sec:rule_join}

Pravilo \texttt{JoinSameClassesRule} (modul
\code{rule_join_same_classes.py}) grupiše časove istog predmeta,
grupe i tipa (predavanje/vežbe) u uzastopne blokove. Cilj je izbeći
rascepkane rasporede gde su, na primer, četiri časa istog predmeta
razbacana kroz nedelju.

\subsection{Formiranje blokova}

Sesije se najpre grupišu u \emph{familije} po ključu
(predmet, grupa, tip). Za svaku familiju, broj sesija se deli na blokove
funkcijom \texttt{split\_into\_blocks}:

\begin{equation}
\text{split}(n) =
\begin{cases}
[n] & \text{ako } n \leq 3 \\
[2, 2] & \text{ako } n = 4
\end{cases}
\label{eq:blocks}
\end{equation}

Tako se tročasovni blok drži zajedno, dok se četiri časa dele na dva
dvočasa. Unutar svakog bloka, prva sesija se bira kao \emph{sidro}
(eng. \emph{anchor}), a ostale se vezuju za nju.

\subsection{Tvrda varijanta}

U tvrdoj varijanti, svaka sesija $s_i$ u bloku (sa rednim brojem $i$ u
odnosu na sidro $a$) mora biti istog dana, u istoj učionici i u
uzastopnom satu:

\begin{align}
\text{day}_{s_i} &= \text{day}_a \\
\text{slot}_{s_i} &= \text{slot}_a + i \\
\text{room}_{s_i} &= \text{room}_a \\
\text{slot}_a + |B| - 1 &\leq H - 1
\end{align}

gde je $|B|$ veličina bloka, a poslednji uslov obezbeđuje da blok stane
unutar radnog vremena. U kodu:

\begin{english}
\begin{lstlisting}[language=Python, caption={Tvrdo spajanje bloka oko sidra}]
for i, s in enumerate(block[1:], start=1):
    solver.model.Add(solver.day_var[s]  == solver.day_var[anchor])
    solver.model.Add(solver.slot_var[s] == solver.slot_var[anchor] + i)
    solver.model.Add(solver.room_var[s] == solver.room_var[anchor])
solver.model.Add(solver.slot_var[anchor] + len(block) - 1 <= H - 1)
\end{lstlisting}
\end{english}

\subsection{Meka varijanta}

U mekoj varijanti, ista ograničenja se uvode uslovno, preko literala
$b$ (``blok je spojen''). Ako blok ne može biti spojen, $b$ je netačno
i uvodi se promenljiva prekršaja:

\begin{english}
\begin{lstlisting}[language=Python, caption={Meko spajanje sa promenljivom prekršaja}]
b = solver.model.NewBoolVar(f"joined_{block[0]}_{block[-1]}")
for i, s in enumerate(block[1:], start=1):
    solver.model.Add(solver.day_var[s]  == solver.day_var[anchor]).OnlyEnforceIf(b)
    solver.model.Add(solver.slot_var[s] == solver.slot_var[anchor] + i).OnlyEnforceIf(b)
    solver.model.Add(solver.room_var[s] == solver.room_var[anchor]).OnlyEnforceIf(b)

violation = solver.model.NewBoolVar(f"violation_{block[0]}_{block[-1]}")
solver.model.Add(violation == 1).OnlyEnforceIf(b.Not())
solver.model.Add(violation == 0).OnlyEnforceIf(b)
\end{lstlisting}
\end{english}

Promenljiva \texttt{violation} ulazi u funkciju cilja
\eqref{eq:objective}, pa rešavač spaja blokove kad god je to
moguće, a krši pravilo samo kada je neophodno.

\section{Pravilo: jedna lokacija po danu za grupu}
\label{sec:rule_location}

Pravilo \texttt{SingleLocationInDayForGroupRule} (modul
\code{rule_single_location_in_day_for_group.py}) sprečava da
grupa studenata menja zgradu (lokaciju) tokom jednog radnog dana.
Matematički fakultet izvodi nastavu na više lokacija (Studentski trg,
zgrada Sv. Nikole, Jagićeva), pa bi promena zgrade između dva uzastopna časa
bila nepraktična. Ovo je tvrdo ograničenje.

\subsection{Mehanizam}

Ključ implementacije je dvostruka primena ograničenja
\texttt{AddElement}. Najpre, za svaku sesiju $s$ uvodimo promenljivu
$\text{loc}_s$ koja predstavlja lokaciju dodeljene učionice:

\begin{equation}
\text{loc}_s = \text{locOfClassroom}[\text{room}_s]
\end{equation}

Pošto je $\text{room}_s$ promenljiva, koristimo \texttt{AddElement} da
indeksiramo niz lokacija učionica:

\begin{english}
\begin{lstlisting}[language=Python, caption={Kanalisanje učionice u lokaciju}]
loc[s] = solver.model.NewIntVar(0, max_loc, f"loc_{s}")
solver.model.AddElement(
    solver.room_var[s], self.locations_of_classrooms, loc[s]
)
\end{lstlisting}
\end{english}

Zatim, za svaku grupu $g$ i dan $d$ uvodimo promenljivu
$\text{groupDayLoc}_{g,d}$ -- lokaciju te grupe tog dana. Za svaku
sesiju $s$ grupe $g$ namećemo da njena lokacija bude jednaka lokaciji
grupe na dan kada se sesija održava:

\begin{equation}
\text{loc}_s = \text{groupDayLoc}_{g}[\text{day}_s]
\end{equation}

Ponovo koristimo \texttt{AddElement}, ovog puta indeksirajući niz
dnevnih lokacija promenljivom $\text{day}_s$:

\begin{english}
\begin{lstlisting}[language=Python, caption={Vezivanje lokacije sesije za lokaciju grupe po danu}]
for g, indices in group_session_indices.items():
    day_locs = [group_day_loc[g, d] for d in range(D)]
    for s in indices:
        # loc[s] == group_day_loc[g][day_var[s]]
        solver.model.AddElement(solver.day_var[s], day_locs, loc[s])
\end{lstlisting}
\end{english}

Pošto sve sesije grupe $g$ u danu $d$ moraju imati istu vrednost
$\text{loc}_s$ (jednaku $\text{groupDayLoc}_{g,d}$), rešavač je prinuđen
da ih sve smesti u učionice na istoj lokaciji.

\section{Pravilo: bez procepa u rasporedu}
\label{sec:rule_nogaps}

Pravilo \texttt{NoGapsInScheduleRule} (modul
\code{rule_no_gaps_in_schedule.py}) cilja na kvalitet rasporeda
iz ugla studenata: izbegava \emph{procepe} -- slobodne sate koji se
nalaze između dva zauzeta sata u istom danu za istu grupu. Na primer,
ako grupa ima čas u 8h i u 11h, a 9h i 10h su slobodni, to su dva
procepa.

\subsection{Matrica zauzetosti}

Najpre se za svaku grupu gradi matrica zauzetosti
$\text{busy}_{d,h}$ koja je tačna ako grupa ima neki čas u danu $d$,
satu $h$. Za svaku sesiju i svaki termin $t = d \cdot H + h$ uvodimo
reifikovani literal koji je tačan kada je sesija u tom terminu, a zatim
ih agregiramo logičkim ILI preko \texttt{AddMaxEquality}:

\begin{english}
\begin{lstlisting}[language=Python, caption={Gradnja matrice zauzetosti}]
for s in sess:
    lit = solver.model.NewBoolVar(f"at_{s}_{t}")
    solver.model.Add(solver.flat_time_var[s] == t).OnlyEnforceIf(lit)
    solver.model.Add(solver.flat_time_var[s] != t).OnlyEnforceIf(lit.Not())
    lits.append(lit)
b = solver.model.NewBoolVar(f"busy_{g}_{d}_{h}")
solver.model.AddMaxEquality(b, lits)   # OR preko svih sesija grupe
\end{lstlisting}
\end{english}

\subsection{Detekcija procepa}

Sat $h$ je procep ako je slobodan ($\neg \text{busy}_{d,h}$), a postoji
zauzet sat i pre i posle njega u istom danu. Pomoćni literali
$\text{hasBefore}$ i $\text{hasAfter}$ se računaju kao logičko ILI nad
prefiksom, odnosno sufiksom reda zauzetosti. Sama promenljiva
$\text{gap}$ se reifikuje kao konjunkcija:

\begin{equation}
\text{gap}_{d,h} \Leftrightarrow
\text{hasBefore}_{d,h} \land \text{hasAfter}_{d,h} \land
\neg \text{busy}_{d,h}
\end{equation}

\begin{english}
\begin{lstlisting}[language=Python, caption={Detekcija unutrašnjeg procepa}]
has_before = solver.model.NewBoolVar(f"hb_{g}_{d}_{h}")
solver.model.AddMaxEquality(has_before, busy_day[:h])
has_after = solver.model.NewBoolVar(f"ha_{g}_{d}_{h}")
solver.model.AddMaxEquality(has_after, busy_day[h + 1:])

gap = solver.model.NewBoolVar(f"gap_{g}_{d}_{h}")
solver.model.AddBoolAnd(
    [has_before, has_after, busy_day[h].Not()]
).OnlyEnforceIf(gap)
solver.model.AddBoolOr(
    [has_before.Not(), has_after.Not(), busy_day[h]]
).OnlyEnforceIf(gap.Not())
\end{lstlisting}
\end{english}

Rubni sati (prvi i poslednji u danu) ne mogu biti unutrašnji procepi,
pa se preskaču.

\subsection{Tvrda i meka varijanta}

U tvrdoj varijanti svaki procep se zabranjuje ograničenjem
$\text{gap}_{d,h} = 0$. U mekoj varijanti, promenljive $\text{gap}_{d,h}$
se vraćaju kao prekršaji i ulaze u funkciju cilja, pa rešavač
minimizuje ukupan broj procepa:

\begin{english}
\begin{lstlisting}[language=Python, caption={Tvrda naspram meke varijante}]
if self.is_hard:
    for gap in gap_vars:
        solver.model.Add(gap == 0)   # zabrana svih procepa
    return []
return gap_vars                       # meko: minimizuj broj procepa
\end{lstlisting}
\end{english}

Ispravnost pravila proverava se testom
\code{test_no_gaps_in_schedule.py}, koji potvrđuje da rešenje na
malom ulazu nema nijedan unutrašnji procep.

\section{Pravilo: kapacitet učionice kao meko ograničenje}
\label{sec:kapacitet}

Svaka učionica ima kapacitet, a svaka sesija poznat broj slušalaca
$\text{size}_s$ (zbir grupa koje je pohađaju, odeljak
\ref{sec:cohorts}). Prirodno je zahtevati da svi slušaoci stanu u
učionicu, $\text{size}_s \leq \text{capacity}_{r_s}$, i model je u
ranijoj verziji to nametao kao tvrd uslov -- učionice nedovoljnog
kapaciteta prosto nisu ulazile u listu dozvoljenih učionica sesije
(odeljak \ref{sec:cp_form}).

Takav tvrd uslov se, međutim, pokazao i previše krut i previše
osetljiv na kvalitet ulaznih podataka:

\begin{itemize}
\item brojevi upisanih studenata su procene, a stvarna prisutnost je
  po pravilu manja od upisa; nerealno je da jedan student preko
  kapaciteta učini ceo raspored nedopustivim,
\item u kombinaciji sa zahtevom za računarima kapacitet drastično sužava
  izbor učionica (računarskih učionica je malo i one su male), pa
  rešavač lako prijavi \texttt{INFEASIBLE} i ne vrati \emph{nikakav}
  raspored -- a raspored sa nekoliko prepunjenih učionica u praksi je
  daleko korisniji od nijednog,
\item svako sužavanje dopustivog skupa gubi informaciju: umesto poruke
  ,,nema rešenja`` korisnije je videti \emph{gde} i \emph{koliko}
  raspored odstupa od željenog.
\end{itemize}

Zbog toga kapacitet prelazi iz tvrdih ograničenja u pravilo
\code{rule_room_capacity.py} (\texttt{roomCapacity}). U mekoj varijanti
pravilo ne sužava dopustiv skup -- sesija sme u bilo koju učionicu koja
zadovoljava tvrda ograničenja -- ali plaća cenu jednaku broju studenata
koji u toj učionici ostaju bez mesta:

\begin{equation}
\text{cost}_s = \max
  \left( 0,\; \text{size}_s - \text{capacity}_{r_s} \right)
\label{eq:room_cost}
\end{equation}

Bitno je da se kažnjava \emph{samo} prekoračenje. Sve učionice u koje
cela kohorta staje su međusobno jednako dobre, pa pravilo ne pravi
nepotrebnu razliku između njih i ne troši prostor rešenja na to da
najveću salu dobije najveća grupa kad to nikome ne pomaže. Kazna je nula
uvek kada raspored nikoga ne ostavlja bez mesta, što je i uobičajen
ishod: prekoračenje se javlja samo tamo gde učionica odgovarajuće
veličine (ili odgovarajuće opremljena računarima) više nije slobodna.

Za razliku od procepa, prekršaj ovde nije ,,jedan događaj`` nego
\emph{količina} -- pravilo je prvo koje kao prekršaj vraća celobrojne, a
ne binarne promenljive. Funkcija cilja \eqref{eq:objective} je težinska
suma vraćenih promenljivih, pa se takva cena uklapa bez izmena u jezgru
rešavača. Cena po prekršaju je pri tome namerno niska, a
niska težina drži ovaj sabirak podređenim ostalim ciljevima čak i kada
se prekoračenje meri desetinama studenata.

U implementaciji se za svaku sesiju gradi tabela prekoračenja po
učionicama (konstante, jer su i veličina sesije i kapaciteti poznati
unapred), a promenljiva prekoračenja se vezuje za izabranu učionicu
preko \texttt{AddElement} (odeljak \ref{sec:cpsat_tehnike}). Sesije koje
staju u svaku učionicu ne dobijaju nikakvu promenljivu:

\begin{english}
\begin{lstlisting}[language=Python, caption={Prekoračenje kapaciteta po sesiji}]
overflows = [max(0, session.size - cap) for cap in capacities]
if not any(overflows):
    continue                      # sesija staje u svaku ucionicu

if self.is_hard:                  # tvrdo: samo ucionice u koje staje
    fitting = [i for i, over in enumerate(overflows) if over == 0]
    solver.model.AddAllowedAssignments(
        [solver.room_var[s]], [[i] for i in fitting])
    continue

overflow = solver.model.NewIntVar(0, max(overflows), f"room_overflow_{s}")
solver.model.AddElement(solver.room_var[s], overflows, overflow)
overflow_vars.append(overflow)
\end{lstlisting}
\end{english}

Tvrda varijanta (\texttt{penalty} 0) svodi se na klasično ograničenje
kapaciteta: izbor učionica se sužava na one u koje kohorta staje, isto
kao što se za računare radi u jezgru rešavača. Ako za neku sesiju takve
učionice nema, greška se prijavljuje odmah pri konstrukciji modela,
umesto da se pojavi kao neinformativan status \texttt{INFEASIBLE}.

Ispravnost obe varijante proverava test \code{test_room_capacity.py}:
meko pravilo bira učionicu u koju grupa staje, a kada takve učionice
nema ipak vraća raspored sa najmanjim prekoračenjem; tvrdo pravilo takve
učionice zabranjuje.


% ------------------------------------------------------------------------------
\chapter{Raspoređivanje nastavnog osoblja}
\label{chp:osoblje}
\chaptermark{Raspoređivanje osoblja}
% ------------------------------------------------------------------------------

Do sada je model raspoređivao samo grupe studenata u učionice i termine.
U ovoj glavi model proširujemo nastavnim osobljem: svakoj sesiji se
dodeljuje nastavnik, uvode se ograničenja vezana za nastavnike, i uvodi
se pojam \emph{kohorte} -- skupa grupa koje zajedno slušaju isti čas.
Ovaj deo predstavlja poslednji korak u približavanju modela realnom
scenariju na Matematičkom fakultetu.

Nastavno osoblje se zadaje u \emph{odvojenom} ulaznom fajlu, tako da
osnovni ulaz (predmeti, učionice, broj upisanih) ostaje nepromenjen.
Time se podaci o osoblju mogu menjati nezavisno, a raspoređivanje bez
osoblja i dalje radi kao ranije. Pokazaće se da je taj fajl i prirodno
mesto za definisanje grupa studenata (odeljak
\ref{sec:grupe_iz_osoblja}): grupa postoji tek onda kada neko drži
nastavu za nju.

\section{Model osoblja i dodela nastave}
\label{sec:staff_model}

Podaci o osoblju modelovani su klasom \texttt{StaffInput} (modul
\code{model.py}) koja sadrži tri celine: listu nastavnika
(\texttt{teachers}), listu dodela nastave (\texttt{assignments}) i
konfiguraciju pravila (\texttt{rules}). Nastavnik
(\texttt{Teacher}) ima identifikator, ime i opciono lično ograničenje
najvećeg broja radnih dana:

\begin{english}
\begin{lstlisting}[language=Python, caption={Modeli nastavnika i dodele nastave}]
@dataclass(config=_config)
class Teacher:
    id: int
    name: str
    max_working_days: Optional[int] = Field(default=None, alias="maxWorkingDays")

@dataclass(config=_config)
class TeachingAssignment:
    teacher_id: int = Field(alias="teacherId")
    course_id: int = Field(alias="courseId")
    session_type: str = Field(alias="sessionType")   # "theory" | "practice"
    group_ids: Optional[List[str]] = Field(default=None, alias="groupIds")
\end{lstlisting}
\end{english}

Dodela (\texttt{TeachingAssignment}) povezuje par (predmet, tip nastave)
sa nastavnikom, i opciono navodi konkretne grupe. Prilikom generisanja
sesija, svaka sesija dobija polje \texttt{teacher\_id}, a rešavač gradi
preslikavanje \texttt{teacher\_sessions} (nastavnik $\rightarrow$ indeksi
njegovih sesija). Sesije bez dodeljenog nastavnika (\texttt{teacher\_id}
jednako \texttt{None}) ne učestvuju u pravilima osoblja.

\subsection{Nastavnik ne može držati dva časa istovremeno}

Prvo i osnovno ograničenje osoblja je tvrdo i dodaje se u jezgru
rešavača: nijedan nastavnik ne može držati dve sesije u istom terminu.
Analogno ograničenju za grupe (HC2, listing \ref{lst:cp_alldiff}),
koristimo \texttt{AllDifferent} nad apsolutnim vremenima svih sesija
jednog nastavnika:

\begin{equation}
\texttt{AllDifferent}(\{\text{flatTime}_s : s \in \mathcal{S}_t\}),
  \quad \forall t
\label{eq:staff_nodouble}
\end{equation}

gde je $\mathcal{S}_t$ skup sesija nastavnika $t$. U kodu
(modul \code{cp_solver.py}):

\begin{english}
\begin{lstlisting}[language=Python, caption={Nastavnik nije dvostruko zauzet}]
for teacher_id, session_indices in self.teacher_sessions.items():
    if len(session_indices) > 1:
        self.model.AddAllDifferent(
            [self.flat_time_var[s] for s in session_indices]
        )
\end{lstlisting}
\end{english}

\section{Grupe studenata izvedene iz dodela nastave}
\label{sec:grupe_iz_osoblja}

Dodela nastave imenuje grupe kojima nastavnik drži čas, na primer
\texttt{"groupIds": ["1i1a", "1i1b", "1i1v"]}. Time je u fajlu sa
osobljem već zapisano \emph{koje} grupe postoje, pa bi njihovo ponovno
navođenje u osnovnom ulazu bilo redudantno.

Postupak je implementiran u funkciji \texttt{build\_groups} (modul
\code{data.py}) i odvija se u dva koraka.

\paragraph{Imena i pripadnost.} Prolazi se kroz sve dodele koje navode
\texttt{groupIds}; za svaki naziv grupe smer i semestar se \emph{ne}
zadaju posebno, nego se čitaju sa predmeta na koji se dodela odnosi
(\texttt{courseId}). Grupa \texttt{1i1a} tako pripada smeru i semestru
predmeta koji sluša. Ako se isti naziv grupe pojavi na predmetima iz
dva različita para (smer, semestar), to je nekonzistentan ulaz i
prijavljuje se kao greška:

\begin{english}
\begin{lstlisting}[language=Python, caption={Imena grupa iz dodela nastave}]
for a in staff_input.assignments:
    if not a.group_ids:
        continue
    course = courses_by_id.get(a.course_id)
    key = (course.track_id, course.semester)
    for gid in a.group_ids:
        existing = meta.get(gid)
        if existing is not None and existing != key:
            raise ValueError(f"Group '{gid}' is used for both {existing} and {key}")
        meta[gid] = key
\end{lstlisting}
\end{english}

\paragraph{Veličine.} Broj studenata se i dalje uzima iz sekcije
\texttt{studentsEnrolled} osnovnog ulaza, koja je zadata po paru
(smer, semestar). Upis se ravnomerno deli na sve grupe tog para, a
ostatak se raspoređuje na prve grupe po abecednom redu naziva, tako da
se veličine razlikuju najviše za jednog studenta:

\begin{english}
\begin{lstlisting}[language=Python, caption={Deljenje upisanih studenata na imenovane grupe}]
base, rem = divmod(enrollment.count, len(group_ids))
for i, gid in enumerate(sorted(group_ids)):
    count = base + (1 if i < rem else 0)
    groups.append(Group(gid, enrollment.track_id, count, enrollment.semester))
\end{lstlisting}
\end{english}

Na primer, 100 upisanih studenata prvog semestra Informatike i šest
grupa navedenih u dodelama (\texttt{1i1a}, \texttt{1i1b},
\texttt{1i1v}, \texttt{1i2a}, \texttt{1i2b}, \texttt{1i2v}) daju četiri
grupe od 17 i dve od 16 studenata.

Time je zbir veličina grupa po konstrukciji jednak broju upisanih: nad
punim letnjim semestrom 44 grupe zbrajaju tačno 919 studenata. To je i
razlika u odnosu na deljenje po veličini iz odeljka \ref{sec:sesije},
gde se ostatak pri deljenju gubi. Fajl sa osobljem, dakle, određuje
\emph{koje} grupe postoje i \emph{ko} im drži nastavu, a osnovni ulaz
\emph{koliko} ih je i \emph{šta} slušaju. Ako neki par (smer, semestar)
nema ni jednu imenovanu grupu u dodelama, za njega se koristi rezervni
postupak deljenja po veličini grupe iz odeljka \ref{sec:sesije}, pa
delimično popunjen fajl sa osobljem ne obara raspoređivanje ostatka
fakulteta.

\section{Kohorte i zajedničke sesije}
\label{sec:cohorts}

Bez dodatne pažnje, ako dve grupe slušaju isto predavanje, model bi
napravio dve odvojene sesije i time primorao nastavnika da isti čas drži
dva puta. Da bismo to izbegli, uvodimo pojam \emph{kohorte}
(\texttt{Cohort}) -- skupa grupa koje zajedno pohađaju sesije jednog
kursa i tipa, sa (opcionim) zajedničkim nastavnikom. Veličina kohorte je
zbir broja studenata svih njenih grupa, što je bitno za kapacitet
učionice.

Kohorte se formiraju funkcijom \texttt{build\_cohorts} (modul
\code{data.py}) na osnovu dodela nastave, po sledećim pravilima:

\begin{itemize}
\item \textbf{Zajednička (joint) dodela} -- ako dodela navodi više grupa
  u \texttt{groupIds}, te grupe čine \emph{jednu} kohortu i dobijaju
  \emph{jednu zajedničku sesiju}. Tako nastavnik drži čas \emph{jednom},
  a prisustvuju mu sve grupe (npr. oba toka informatike slušaju zajedno
  predavanje iz predmeta ``Programske paradigme'').
\item \textbf{Dodela za jednu grupu} -- \texttt{groupIds} sa jednom
  grupom važi samo za tu grupu.
\item \textbf{Generička dodela} -- ako \texttt{groupIds} nije naveden,
  dodela važi za svaku \emph{preostalu} grupu kursa pojedinačno (svaka
  grupa dobija svoju sesiju sa tim nastavnikom).
\item Grupe bez ijedne dodele postaju pojedinačne kohorte bez nastavnika.
\end{itemize}

\begin{english}
\begin{lstlisting}[language=Python, caption={Formiranje kohorti iz dodela nastave}]
for a in assignments:
    if a.course_id != course.id or a.session_type != session_type:
        continue
    if a.group_ids is None:
        generic_teacher = a.teacher_id          # primenjuje se na sve preostale grupe
        continue
    cohorts.append(Cohort([groups_by_id[g] for g in a.group_ids], a.teacher_id))
    for g in a.group_ids:
        remaining.pop(g, None)                  # ove grupe su pokrivene

for g in remaining.values():                    # grupe bez eksplicitne dodele
    cohorts.append(Cohort([g], generic_teacher))
\end{lstlisting}
\end{english}

Za svaku kohortu i svaki traženi čas iz kvote kreira se tačno jedna
sesija, čija je lista grupa \texttt{group\_ids} jednaka grupama kohorte.
Zajednička sesija tako ulazi u ograničenje o koliziji \emph{svake} grupe
koja je pohađa (tvrdo ograničenje HC2 prolazi kroz sve
\texttt{session.group\_ids}). Veličina kohorte ne ograničava izbor
učionice tvrdo, ali kroz meko pravilo iz odeljka \ref{sec:kapacitet}
usmerava sesiju ka dovoljno velikoj sali.

\paragraph{Primer.} Dve grupe \texttt{ga} (30 studenata) i \texttt{gb}
(25 studenata) slušaju predavanja zajedno, a vežbe odvojeno. Za predmet
sa 2 predavanja i 2 vežbe nedeljno dobija se: 2 zajedničke teorijske
sesije (kohorta \texttt{ga+gb}, veličine 55) i po 2 vežbe za svaku grupu
-- ukupno 6 sesija umesto 8. Zajedničkom predavanju je za 55 slušalaca
potrebna sala kapaciteta najmanje 55, dok vežbama od 30, odnosno 25
studenata odgovara i znatno manja učionica.

\section{Pravila nastavnog osoblja}
\label{sec:staff_rules}

Pravila osoblja koriste istu \emph{plugin} arhitekturu i CP-SAT tehnike
(reifikacija, \texttt{AddMaxEquality}, \texttt{AddBoolAnd}/\texttt{AddBoolOr},
\texttt{AddElement}) uvedene u glavi \ref{chp:nadogradnja}. Konfigurišu
se u \texttt{rules} sekciji staff fajla i mogu biti tvrda
(\texttt{penalty} 0) ili meka (\texttt{penalty} > 0).

\subsection{Jedna lokacija po danu za nastavnika}

Pravilo \code{rule_staff_single_location_in_day.py}
(\texttt{staffSingleLocationInDay}) sprečava da nastavnik menja zgradu
tokom jednog dana -- fizički bi bilo neizvodljivo držati dva uzastopna
časa na različitim lokacijama. Mehanizam je isti kao za grupe: za svakog
nastavnika i dan uvodi se promenljiva lokacije, a sve njegove sesije tog
dana vezuju se za nju pomoću \texttt{AddElement}:

\begin{english}
\begin{lstlisting}[language=Python, caption={Nastavnik na jednoj lokaciji dnevno}]
day_locs = [model.NewIntVar(0, max_loc, f"teacher_day_loc_{teacher_id}_{d}")
            for d in range(D)]
for s in session_indices:
    # loc_var[s] == day_locs[day_var[s]]
    solver.model.AddElement(solver.day_var[s], day_locs, solver.loc_var[s])
\end{lstlisting}
\end{english}

Ovo je tvrdo ograničenje i vraća praznu listu prekršaja.

\subsection{Najveći broj radnih dana nastavnika}

Pravilo \code{rule_staff_max_working_days.py}
(\texttt{staffMaxWorkingDays}) ograničava koliko različitih dana u
nedelji nastavnik dolazi na fakultet. Za svaki dan $d$ uvodi se indikator
$\text{dayUsed}_{t,d}$ (logičko ILI reifikovanih jednakosti
$\text{day}_s = d$ preko sesija nastavnika), a broj radnih dana je
njihov zbir:

\begin{equation}
\text{worked}_t = \sum_{d=0}^{D-1} \text{dayUsed}_{t,d}
\end{equation}

Tvrda varijanta nameće $\text{worked}_t \leq \text{limit}_t$, gde je
$\text{limit}_t$ lični limit nastavnika (\texttt{maxWorkingDays}) ili
podrazumevana vrednost pravila (\texttt{maxDays}). Meka varijanta uvodi
celobrojnu promenljivu viška i minimizuje ga:

\begin{equation}
\min \sum_{t} \text{penalty} \cdot \max(0,\ \text{worked}_t - \text{limit}_t)
\label{eq:staff_days}
\end{equation}

\begin{english}
\begin{lstlisting}[language=Python, caption={Ograničenje broja radnih dana}]
worked_days = sum(day_used)                 # day_used[d] = OR preko sesija
if self.is_hard:
    solver.model.Add(worked_days <= limit)
else:
    excess = solver.model.NewIntVar(0, D, f"staff_days_excess_{teacher_id}")
    solver.model.Add(excess >= worked_days - limit)
    violations.append(excess)               # ulazi u funkciju cilja
\end{lstlisting}
\end{english}

\subsection{Nedeljni budžet procepa (broj sati pauze) nastavnika}

Pravilo \code{rule_staff_no_gap_greater_than.py}
(\texttt{staffMaxGapHoursPerWeek}) ograničava ukupan broj procepa
(pauza) nastavnika \emph{tokom cele nedelje}. Procep se detektuje isto kao u
pravilu \nameref{sec:rule_nogaps}: matrica zauzetosti se gradi
reifikovanim jednakostima i \texttt{AddMaxEquality}, a procep je
konjunkcija ``ima čas pre'', ``ima čas posle'' i ``slobodan sat''. Za
razliku od pravila za grupe, ovde se procepi \emph{sabiraju kroz sve
dane}:

\begin{equation}
\text{weekGaps}_t = \sum_{d=0}^{D-1} \sum_{h} \text{gap}_{t,d,h}
\end{equation}

Tvrda varijanta nameće $\text{weekGaps}_t \leq B$ (budžet
$B = \texttt{maxGapHours}$, podrazumevano 1). Meka varijanta minimizuje
prekoračenje budžeta:

\begin{equation}
\min \sum_{t} \text{penalty} \cdot \max(0,\ \text{weekGaps}_t - B)
\label{eq:staff_gaps}
\end{equation}

\begin{english}
\begin{lstlisting}[language=Python, caption={Nedeljni budžet procepa}]
total_gaps = sum(week_gaps)
if self.is_hard:
    solver.model.Add(total_gaps <= self.max_gap_hours)
else:
    excess = solver.model.NewIntVar(0, D * H, f"staff_gap_excess_{teacher_id}")
    solver.model.Add(excess >= total_gaps - self.max_gap_hours)
    violations.append(excess)
\end{lstlisting}
\end{english}

Svi prekršaji mekih pravila osoblja ulaze u istu globalnu funkciju cilja
iz glave \ref{chp:nadogradnja}:
$\min \sum_{r} \text{penalty}_r \cdot \sum_{v \in V_r} v$.

\section{Konfiguracija i provera}
\label{sec:staff_config}

Osoblje se zadaje u zasebnom JSON fajlu (na primer
\code{staff_2_semester.json}). Sekcija \texttt{rules} bira koja pravila
su aktivna i sa kojom cenom, \texttt{teachers} navodi nastavnike, a
\texttt{assignments} povezuje predmete sa nastavnicima i istovremeno
definiše grupe studenata (odeljak \ref{sec:grupe_iz_osoblja}):

\begin{english}
\begin{lstlisting}[language=Python, caption={Isečak staff ulaznog fajla}]
{
  "rules": {
    "staffMaxWorkingDays":     {"enabled": true, "penalty": 3, "params": {"maxDays": 4}},
    "staffMaxGapHoursPerWeek": {"enabled": true, "penalty": 0, "params": {"maxGapHours": 1}},
    "staffSingleLocationInDay":{"enabled": true, "penalty": 0}
  },
  "teachers": [
    {"id": 4, "name": "milan.mitreski", "maxWorkingDays": 2}
  ],
  "assignments": [
    {"teacherId": 43, "courseId": 54, "sessionType": "theory", "groupIds": ["2la", "2lb"]}
  ]
}
\end{lstlisting}
\end{english}

Rešavač se pokreće sa dodatnim argumentom \texttt{-{}-staff}:

\begin{english}
\begin{verbatim}
bazel run //src/algo:run_cp_solver -- \
    --input input_full_2_semester.json \
    --staff staff_2_semester.json --time-limit 600
\end{verbatim}
\end{english}

Ispravnost pravila osoblja proverava se testovima u
\code{test_staff_rules.py}: da nastavnik nije dvostruko zauzet, da se
poštuje najveći broj radnih dana (uključujući lične limite), da je
nedeljni budžet procepa ispunjen, da nastavnik ne menja lokaciju u toku
dana, i da se zajednička sesija dve grupe rešava kao jedan čas bez
kolizija u rasporedu obe grupe. Izvođenje grupa iz dodela pokriveno je
testovima u \code{test_data.py}.


% ------------------------------------------------------------------------------
\chapter{Završna evaluacija proširenog modela}
\label{chp:finalna-evaluacija}
\chaptermark{Završna evaluacija}
% ------------------------------------------------------------------------------

Evaluacija iz glave \ref{chp:evaluacija} poredila je dve paradigme na
osnovnom modelu, u režimu traženja bilo kog dopustivog rasporeda. Pošto
je model u glavama \ref{chp:nadogradnja} i \ref{chp:osoblje} proširen
mekim ograničenjima, funkcijom cilja i raspoređivanjem nastavnog osoblja,
u ovoj glavi merimo kako se tako prošireni CP-SAT model ponaša na tri
skale problema. Za razliku od prve evaluacije, ovde nas ne zanima samo
da li rešenje postoji, nego i \emph{kvalitet} rešenja i brzina kojom se
on popravlja, pa pratimo vrednost funkcije cilja u zavisnosti od
proteklog vremena izračunavanja.

Merenja su sprovedena u dva okruženja. Prvo je realno okruženje, sa
podacima onakvim kakvi jesu, i ono odgovara na pitanje da li je
prošireni model upotrebljiv u praksi. Drugo je zaoštreno okruženje
(odeljak \ref{sec:finalna_zaostreno}), u kome je broj raspoloživih
termina smanjen tako da meka ograničenja dođu u stvarni sukob, i ono
pokazuje kako se rešavač ponaša na granici rešivosti.

\section{Metodologija}
\label{sec:finalna_metodologija}

\subsection{Instance po godinama studija}

Prošireni model radi nad realnim podacima letnjeg semestra
(\code{input_full_2_semester.json}) i dodelama nastave iz
\code{staff_2_semester.json}, iz kojih se izvode grupe studenata
(odeljak \ref{sec:grupe_iz_osoblja}). Svaka grupa nasleđuje semestar
predmeta koji sluša (2, 4, 6 ili 8), što odgovara godini studija, te se
instanca može seći po godinama:

\begin{enumerate}
\item \textbf{FINAL-S} -- 1. godina (semestar 2),
\item \textbf{FINAL-M} -- 1. i 2. godina (semestri 2 i 4),
\item \textbf{FINAL-L} -- sve četiri godine (semestri 2, 4, 6 i 8).
\end{enumerate}

Skalira se samo potražnja: sve tri instance zadržavaju svih 29 učionica
i sve tri lokacije, pa se razlikuju jedino po količini nastave koju
treba smestiti. Filtriraju se samo predmeti i upisi; grupe se povlače
same, jer se izvode iz dodela za predmete koji su ostali, a dodele za
izbačene predmete se prosto ne aktiviraju -- ulaz sa osobljem tako
ostaje nepromenjen.

Ovako definisane instance poklapaju se sa podskupovima iz glave
\ref{chp:evaluacija}: isti je ulazni fajl, iste godine studija i isti
skup učionica, tako da se rezultati pre i posle proširenja modela
odnose na istu nastavu. U konfiguraciji bez osoblja skup sesija je
doslovno isti (204, 398 i 722 sesije po skalama). Kada se uključi
osoblje, grupe se izvode iz dodela nastave umesto deljenjem upisa na
grupe od najviše 50 studenata, pa se broj sesija menja -- raste broj
imenovanih grupa, a predavanja koja više grupa sluša zajedno spajaju se
u jednu sesiju.

\subsection{Konfiguracije i funkcija cilja}

Svaka skala se rešava u dve konfiguracije:

\begin{itemize}
\item \textbf{bez osoblja} -- pravila iz glave \ref{chp:nadogradnja}:
  spajanje časova u blokove kao tvrdo ograničenje, izbegavanje procepa
  kao meko ograničenje sa cenom 2 i kapacitet učionica kao meko
  ograničenje sa cenom 1;
\item \textbf{sa osobljem} -- dodatno i ceo model iz glave
  \ref{chp:osoblje}: kohorte i zajedničke sesije, tvrdo ograničenje da
  nastavnik nije dvostruko zauzet, tvrda pravila za jednu lokaciju po
  danu i nedeljni budžet procepa nastavnika, i meko pravilo najvećeg
  broja radnih dana sa cenom 3.
\end{itemize}

Funkcija cilja je, u skladu sa izrazom \eqref{eq:objective} iz odeljka
\ref{sec:arhitektura_pravila}, težinska suma prekršaja mekih pravila:
\[
  \min \; 2 \cdot \underbrace{\sum_{g,d,h} \text{gap}_{g,d,h}}_{\text{procepi grupa}}
  \;+\; 3 \cdot \underbrace{\sum_{t} \text{excess}_{t}}_{\text{dani preko limita nastavnika}}
  \;+\; 1 \cdot \underbrace{\sum_{s} \text{overflow}_{s}}_{\text{studenti bez mesta}},
\]
pri čemu u konfiguraciji bez osoblja drugi sabirak ne postoji. Treći
sabirak je meko ograničenje kapaciteta iz odeljka \ref{sec:kapacitet} i
u obe konfiguracije nosi najmanju težinu. Vrednost 0 znači da nijedno
meko ograničenje nije prekršeno, dakle raspored bez ijednog procepa, bez
ijednog nastavnika koji radi više dana nego što mu je dozvoljeno i bez
ijednog studenta koji ostaje bez mesta u učionici.

\subsection{Uslovi merenja}

Merenja su izvršena na računaru sa procesorom Apple M4 (10 jezgara) i
24\,GB radne memorije, pod operativnim sistemom macOS 26.5, sa
bibliotekom OR-Tools 9.15. Rešavač koristi 8 niti pretrage i vremenski
limit od 600 sekundi po pokretanju.

Vrednost funkcije cilja kroz vreme prati se povratnim pozivom
(\emph{callback}) koji CP-SAT poziva pri svakom pronađenom poboljšanju.
Za svako poboljšanje beleže se proteklo vreme pretrage, vrednost cilja i
najbolja donja granica, iz čega se dobija stepenasta funkcija: vrednost
u trenutku $t$ je cilj poslednjeg rešenja nađenog do tog trenutka.
Rešenja se, kao i u prvoj evaluaciji, nezavisno proveravaju funkcijom
\code{validate_solution}, koja je za potrebe proširenog modela dopunjena
i proverom da nastavnik nema dva časa u istom terminu.

\section{Veličina modela}
\label{sec:finalna_velicina}

Tabela \ref{tbl:final_size} prikazuje veličinu modela po skalama i
konfiguracijama.

\begin{table}[ht]
\centering
\caption{Veličina proširenog CP modela po skalama}
\label{tbl:final_size}
\begin{tabular}{lrrrrrr}
\toprule
& \multicolumn{2}{c}{\textbf{FINAL-S}}
& \multicolumn{2}{c}{\textbf{FINAL-M}}
& \multicolumn{2}{c}{\textbf{FINAL-L}} \\
\cmidrule(lr){2-3} \cmidrule(lr){4-5} \cmidrule(lr){6-7}
\textbf{Metrika} & bez & sa & bez & sa & bez & sa \\
\midrule
Sesije             & $204$      & $216$      & $398$      & $441$       & $722$      & $839$ \\
Nastavnici         & $0$        & $39$       & $0$        & $81$        & $0$        & $112$ \\
Promenljive        & $15\,348$  & $41\,036$  & $30\,026$  & $88\,223$   & $54\,652$  & $144\,685$ \\
Ograničenja        & $27\,857$  & $72\,082$  & $54\,424$  & $155\,803$  & $98\,974$  & $257\,191$ \\
Vreme konstrukcije & 0,09 s     & 0,24 s     & 0,18 s     & 0,52 s      & 0,35 s     & 0,94 s \\
\bottomrule
\end{tabular}
\end{table}

Kolone označene sa ,,bez`` odnose se na konfiguraciju bez osoblja, a
kolone ,,sa`` na konfiguraciju sa raspoređivanjem nastavnog osoblja.

Tabela \ref{tbl:final_tok} sumira tok rešavanja: kada je nađeno prvo
rešenje, kolika je bila njegova vrednost cilja, kada je dostignut
najbolji rezultat i koliko je poboljšanja usput pronađeno.

\begin{table}[ht]
\centering
\caption{Tok rešavanja proširenog modela}
\label{tbl:final_tok}
\begin{tabular}{lrrrrrr}
\toprule
& \multicolumn{2}{c}{\textbf{FINAL-S}}
& \multicolumn{2}{c}{\textbf{FINAL-M}}
& \multicolumn{2}{c}{\textbf{FINAL-L}} \\
\cmidrule(lr){2-3} \cmidrule(lr){4-5} \cmidrule(lr){6-7}
\textbf{Metrika} & bez & sa & bez & sa & bez & sa \\
\midrule
Prvo rešenje         & 1,59 s  & 2,69 s & 4,86 s   & 9,06 s   & 23,92 s & 30,03 s \\
Cilj prvog rešenja   & $1\,355$ & $585$ & $1\,009$ & $1\,849$ & $2\,440$ & $2\,973$ \\
Broj nađenih rešenja & $23$    & $10$   & $15$     & $38$     & $23$    & $56$ \\
Vreme do najboljeg   & 2,35 s  & 3,79 s & 6,94 s   & 16,37 s  & 40,84 s & 120,44 s \\
Ukupno vreme         & 2,46 s  & 4,05 s & 7,15 s   & 16,95 s  & 41,23 s & 121,76 s \\
Najbolji cilj        & $0$     & $0$    & $0$      & $0$      & $0$     & $0$ \\
Donja granica        & $0$     & $0$    & $0$      & $0$      & $0$     & $0$ \\
Status               & \multicolumn{6}{c}{\textsc{optimal}} \\
Rešenje validno      & \multicolumn{6}{c}{da} \\
\bottomrule
\end{tabular}
\end{table}

\section{Funkcija cilja u zavisnosti od vremena}
\label{sec:finalna_cilj_vreme}

Tabela \ref{tbl:final_checkpoints} prikazuje vrednost funkcije cilja u
unapred zadatim vremenskim presecima. Crtica znači da u tom trenutku
rešavač još nije imao nijedno dopustivo rešenje.

\begin{table}[ht]
\centering
\caption{Vrednost funkcije cilja u vremenskim presecima}
\label{tbl:final_checkpoints}
\begin{tabular}{lrrrrrr}
\toprule
& \multicolumn{2}{c}{\textbf{FINAL-S}}
& \multicolumn{2}{c}{\textbf{FINAL-M}}
& \multicolumn{2}{c}{\textbf{FINAL-L}} \\
\cmidrule(lr){2-3} \cmidrule(lr){4-5} \cmidrule(lr){6-7}
\textbf{Vreme} & bez & sa & bez & sa & bez & sa \\
\midrule
5 s    & $0$ & $0$ & $929$ & --       & --    & --    \\
10 s   & $0$ & $0$ & $0$   & $1\,008$ & --    & --    \\
20 s   & $0$ & $0$ & $0$   & $0$      & --    & --    \\
30 s   & $0$ & $0$ & $0$   & $0$      & $519$ & --    \\
60 s   & $0$ & $0$ & $0$   & $0$      & $0$   & $140$ \\
120 s  & $0$ & $0$ & $0$   & $0$      & $0$   & $2$   \\
300 s  & $0$ & $0$ & $0$   & $0$      & $0$   & $0$   \\
600 s  & $0$ & $0$ & $0$   & $0$      & $0$   & $0$   \\
\bottomrule
\end{tabular}
\end{table}

Pošto se sva poboljšanja dešavaju u uskom intervalu odmah po nalaženju
prvog rešenja, gruba mreža preseka ne pokazuje dinamiku pretrage. Zato
u tabelama \ref{tbl:final_traj_s}, \ref{tbl:final_traj_m} i
\ref{tbl:final_traj_l} dajemo celu putanju poboljšanja za konfiguraciju
sa osobljem. Konfiguracija bez osoblja prolazi kroz istu vrstu putanje
(23, 15 i 23 pronađena rešenja po skalama), samo je kraća i celom
dužinom se odvija u prvih nekoliko sekundi, pa je ne prikazujemo
posebno.

\begin{table}[ht]
\centering
\caption{Putanja poboljšanja, FINAL-S sa osobljem}
\label{tbl:final_traj_s}
\begin{tabular}{rr}
\toprule
\textbf{Vreme} & \textbf{Cilj} \\
\midrule
2,69 s & $585$ \\
2,72 s & $557$ \\
2,75 s & $555$ \\
2,88 s & $384$ \\
2,99 s & $275$ \\
3,24 s & $106$ \\
3,38 s & $48$ \\
3,64 s & $12$ \\
3,66 s & $10$ \\
3,79 s & $0$ \\
\bottomrule
\end{tabular}
\end{table}

\begin{table}[ht]
\centering
\caption{Putanja poboljšanja, FINAL-M sa osobljem}
\label{tbl:final_traj_m}
\begin{tabular}{rr@{\hskip 2.5em}rr}
\toprule
\textbf{Vreme} & \textbf{Cilj} & \textbf{Vreme} & \textbf{Cilj} \\
\midrule
9,06 s  & $1\,849$ & 14,87 s & $32$ \\
9,18 s  & $1\,789$ & 14,95 s & $30$ \\
9,29 s  & $1\,563$ & 15,04 s & $29$ \\
9,89 s  & $1\,008$ & 15,10 s & $27$ \\
10,09 s & $918$    & 15,18 s & $26$ \\
10,16 s & $792$    & 15,29 s & $24$ \\
10,48 s & $667$    & 15,37 s & $23$ \\
10,64 s & $501$    & 15,44 s & $21$ \\
10,66 s & $453$    & 15,51 s & $20$ \\
11,30 s & $346$    & 15,64 s & $18$ \\
11,52 s & $262$    & 15,70 s & $15$ \\
11,74 s & $252$    & 15,77 s & $12$ \\
12,40 s & $170$    & 15,84 s & $11$ \\
13,29 s & $136$    & 15,90 s & $9$ \\
13,71 s & $60$     & 15,99 s & $8$ \\
13,89 s & $57$     & 16,10 s & $6$ \\
14,61 s & $53$     & 16,22 s & $5$ \\
14,74 s & $35$     & 16,28 s & $3$ \\
14,80 s & $33$     & 16,37 s & $0$ \\
\bottomrule
\end{tabular}
\end{table}

\begin{table}[ht]
\centering
\caption{Putanja poboljšanja, FINAL-L sa osobljem}
\label{tbl:final_traj_l}
\begin{tabular}{rr@{\hskip 2em}rr@{\hskip 2em}rr}
\toprule
\textbf{Vreme} & \textbf{Cilj} & \textbf{Vreme} & \textbf{Cilj}
& \textbf{Vreme} & \textbf{Cilj} \\
\midrule
30,03 s & $2\,973$ & 47,18 s & $472$ & 75,22 s  & $59$ \\
30,27 s & $2\,799$ & 47,76 s & $404$ & 77,56 s  & $49$ \\
30,49 s & $2\,612$ & 50,56 s & $386$ & 78,91 s  & $46$ \\
30,94 s & $2\,332$ & 51,01 s & $265$ & 80,18 s  & $44$ \\
31,13 s & $2\,322$ & 55,60 s & $241$ & 82,52 s  & $42$ \\
31,37 s & $2\,319$ & 55,80 s & $179$ & 84,44 s  & $40$ \\
31,39 s & $2\,303$ & 56,63 s & $157$ & 88,20 s  & $25$ \\
31,66 s & $2\,259$ & 58,72 s & $153$ & 90,92 s  & $23$ \\
32,06 s & $1\,823$ & 59,17 s & $152$ & 94,97 s  & $21$ \\
34,00 s & $1\,527$ & 59,67 s & $140$ & 97,65 s  & $15$ \\
35,76 s & $1\,173$ & 63,10 s & $138$ & 103,78 s & $12$ \\
37,81 s & $1\,070$ & 64,72 s & $87$  & 108,77 s & $11$ \\
38,54 s & $750$    & 66,50 s & $80$  & 109,72 s & $9$ \\
39,74 s & $746$    & 67,32 s & $78$  & 111,12 s & $8$ \\
40,10 s & $730$    & 67,55 s & $72$  & 112,46 s & $5$ \\
40,52 s & $709$    & 67,59 s & $70$  & 118,84 s & $3$ \\
41,23 s & $691$    & 67,81 s & $69$  & 119,81 s & $2$ \\
44,85 s & $647$    & 68,51 s & $67$  & 120,44 s & $0$ \\
46,57 s & $533$    & 71,28 s & $65$  &          &      \\
\bottomrule
\end{tabular}
\end{table}

\textbf{Napomena:} prikazani rezultati dobijeni su komandom:

\begin{english}
\begin{verbatim}
bazel run //src/algo:benchmark -- --mode final --max-time 600 \
    --final-json report_final.json
\end{verbatim}
\end{english}

\section{Zaoštreno okruženje}
\label{sec:finalna_zaostreno}

Rezultati iz prethodnog odeljka pokazuju da rešavač na svim skalama
dostigne cilj 0, i to za nekoliko desetina sekundi, najviše dva minuta.
Razlog je što realna instanca ima dovoljno slobodnog prostora: 29
učionica puta 60 termina nedeljno daje 1740 mesta za najviše 839 časova,
dakle popunjenost ispod 50\%. U tako labavom rasporedu meka ograničenja
retko dolaze u stvarni sukob, pa kriva funkcije cilja sigurno padne na
nulu i o granicama modela govori malo.

Da bismo videli kako se model ponaša kada resursi postanu oskudni,
uvodimo \emph{zaoštreno okruženje} -- istu instancu sa dve izmene:

\begin{itemize}
\item nastava traje od 8 do 15 časova umesto do 20, čime broj termina
  pada sa 1740 na 1015, a popunjenost najveće instance raste na 83\%;
\item podrazumevani najveći broj radnih dana nastavnika smanjen je sa 4
  na 2 (lični limiti pojedinih nastavnika, koji su stroži,
  ostaju nepromenjeni).
\end{itemize}

Dužina dana je birana tako da problem ostane jedva rešiv. Skraćenje
za još jedan čas (8--14, dakle 870 termina i 96\% popunjenosti) dovodi
najveću instancu sa osobljem preko granice rešivosti: rešavač za 600
sekundi ne nađe \emph{nijedno} dopustivo rešenje i vraća status
\textsc{unknown}. Zona u kojoj je ovaj problem istovremeno rešiv i
netrivijalan je, dakle, uska -- meri se jednim časom nastave dnevno.

Tabela \ref{tbl:tight_tok} prikazuje tok rešavanja u tako zaoštrenom
okruženju.

\begin{table}[ht]
\centering
\caption{Tok rešavanja u zaoštrenom okruženju (8--15 časova)}
\label{tbl:tight_tok}
\begin{tabular}{lrrrrrr}
\toprule
& \multicolumn{2}{c}{\textbf{FINAL-S}}
& \multicolumn{2}{c}{\textbf{FINAL-M}}
& \multicolumn{2}{c}{\textbf{FINAL-L}} \\
\cmidrule(lr){2-3} \cmidrule(lr){4-5} \cmidrule(lr){6-7}
\textbf{Metrika} & bez & sa & bez & sa & bez & sa \\
\midrule
Promenljive          & $9\,448$  & $24\,511$ & $18\,476$ & $52\,648$ & $33\,602$ & $86\,685$ \\
Ograničenja          & $16\,657$ & $42\,782$ & $32\,524$ & $92\,453$ & $59\,124$ & $152\,891$ \\
Prvo rešenje         & 1,01 s    & 1,63 s    & 4,01 s    & 7,18 s    & 43,26 s   & 67,06 s \\
Cilj prvog rešenja   & $744$     & $684$     & $1\,837$  & $1\,677$  & $2\,545$  & $2\,456$ \\
Broj nađenih rešenja & $18$      & $16$      & $14$      & $23$      & $24$      & $62$ \\
Vreme do najboljeg   & 1,35 s    & 2,73 s    & 6,51 s    & 14,47 s   & 80,31 s   & 209,53 s \\
Najbolji cilj        & $0$       & $3$       & $0$       & $6$       & $0$       & $6$ \\
Donja granica        & $0$       & $3$       & $0$       & $6$       & $0$       & $6$ \\
Status               & \multicolumn{6}{c}{\textsc{optimal}} \\
Rešenje validno      & \multicolumn{6}{c}{da} \\
\bottomrule
\end{tabular}
\end{table}

Za razliku od realnog okruženja, ovde najbolji cilj \emph{nije} nula:
na najmanjoj skali ostaje 3, a na srednjoj i najvećoj 6. Pošto se donja
granica poklapa sa nađenim rešenjem, reč je o dokazano najboljem
mogućem kompromisu -- u zbijenoj nedelji prosto nije moguće da svi
nastavnici stanu u dva radna dana bez ijednog procepa kod grupa.

Sada i mreža vremenskih preseka ima šta da pokaže. Tabela
\ref{tbl:tight_checkpoints} daje vrednost funkcije cilja u zadatim
trenucima.

\begin{table}[ht]
\centering
\caption{Funkcija cilja u vremenskim presecima, zaoštreno okruženje}
\label{tbl:tight_checkpoints}
\begin{tabular}{lrrrrrr}
\toprule
& \multicolumn{2}{c}{\textbf{FINAL-S}}
& \multicolumn{2}{c}{\textbf{FINAL-M}}
& \multicolumn{2}{c}{\textbf{FINAL-L}} \\
\cmidrule(lr){2-3} \cmidrule(lr){4-5} \cmidrule(lr){6-7}
\textbf{Vreme} & bez & sa & bez & sa & bez & sa \\
\midrule
10 s  & $0$ & $3$ & $0$ & $329$ & --    & --    \\
20 s  & $0$ & $3$ & $0$ & $6$   & --    & --    \\
30 s  & $0$ & $3$ & $0$ & $6$   & --    & --    \\
40 s  & $0$ & $3$ & $0$ & $6$   & --    & --    \\
50 s  & $0$ & $3$ & $0$ & $6$   & $874$ & --    \\
60 s  & $0$ & $3$ & $0$ & $6$   & $289$ & --    \\
90 s  & $0$ & $3$ & $0$ & $6$   & $0$   & $472$ \\
120 s & $0$ & $3$ & $0$ & $6$   & $0$   & $100$ \\
300 s & $0$ & $3$ & $0$ & $6$   & $0$   & $6$   \\
600 s & $0$ & $3$ & $0$ & $6$   & $0$   & $6$   \\
\bottomrule
\end{tabular}
\end{table}

Kolona ,,FINAL-L sa`` je ono što se u ovakvim merenjima obično i očekuje:
prvo rešenje se pojavi tek posle nešto više od minuta, a cilj zatim pada
sa 472 na 100 i konačno na 6, kroz gotovo tri i po minuta pretrage. Cela
putanja poboljšanja za to pokretanje data je u tabeli
\ref{tbl:tight_traj_l}.

\begin{table}[ht]
\centering
\caption{Putanja poboljšanja, FINAL-L sa osobljem u zaoštrenom okruženju}
\label{tbl:tight_traj_l}
\begin{tabular}{rr@{\hskip 2em}rr@{\hskip 2em}rr}
\toprule
\textbf{Vreme} & \textbf{Cilj} & \textbf{Vreme} & \textbf{Cilj}
& \textbf{Vreme} & \textbf{Cilj} \\
\midrule
67,06 s & $2\,456$ & 95,43 s  & $300$ & 131,90 s & $61$ \\
67,22 s & $2\,413$ & 98,57 s  & $281$ & 131,94 s & $58$ \\
71,35 s & $2\,025$ & 99,05 s  & $267$ & 135,01 s & $55$ \\
71,53 s & $1\,978$ & 100,51 s & $227$ & 135,26 s & $54$ \\
72,03 s & $1\,871$ & 102,28 s & $205$ & 141,89 s & $51$ \\
72,42 s & $1\,678$ & 102,30 s & $202$ & 144,99 s & $48$ \\
72,65 s & $1\,646$ & 103,85 s & $171$ & 146,64 s & $45$ \\
73,86 s & $1\,493$ & 106,47 s & $157$ & 148,90 s & $42$ \\
75,46 s & $1\,245$ & 106,99 s & $138$ & 149,57 s & $39$ \\
75,82 s & $1\,232$ & 107,26 s & $121$ & 150,68 s & $36$ \\
78,47 s & $1\,158$ & 108,58 s & $116$ & 151,06 s & $34$ \\
81,15 s & $988$    & 109,10 s & $114$ & 152,24 s & $28$ \\
82,97 s & $845$    & 109,17 s & $110$ & 155,16 s & $25$ \\
82,99 s & $843$    & 110,29 s & $103$ & 156,43 s & $21$ \\
84,55 s & $841$    & 112,17 s & $101$ & 158,18 s & $18$ \\
84,99 s & $689$    & 112,26 s & $100$ & 163,43 s & $15$ \\
85,99 s & $664$    & 120,63 s & $86$  & 164,17 s & $12$ \\
88,57 s & $472$    & 126,71 s & $70$  & 166,54 s & $11$ \\
90,25 s & $470$    & 127,18 s & $67$  & 173,26 s & $9$  \\
92,01 s & $406$    & 130,02 s & $66$  & 209,53 s & $6$  \\
93,31 s & $384$    & 131,89 s & $62$  &          &      \\
\bottomrule
\end{tabular}
\end{table}

Zaoštreno okruženje se dobija dodavanjem opcije \code{--profile tight}:

\begin{english}
\begin{verbatim}
bazel run //src/algo:benchmark -- --mode final --profile tight \
    --max-time 600 --final-json report_final_tight.json
\end{verbatim}
\end{english}

\section{Analiza rezultata}
\label{sec:finalna_analiza}

\subsection{Cena proširenja modela}

Uvođenje nastavnog osoblja model uvećava oko dva i po puta: broj
promenljivih raste 2,6 do 2,9 puta, a broj ograničenja 2,6 do 2,9 puta
na svim skalama. Deo tog rasta dolazi od pomoćnih promenljivih kojima se
izražavaju pravila osoblja, a deo od toga što instanca sa osobljem
opisuje \emph{više} nastave: na najvećoj skali 839 naspram 722 časa.

Razlog za veći broj časova je to što se grupe izvode iz dodela nastave
(odeljak \ref{sec:grupe_iz_osoblja}). Bez fajla sa osobljem upis se deli
na grupe od najviše 50 studenata, pa najveća skala ima 30 grupa; dodele
nastave, međutim, opisuju stvarnu podelu na 44 grupe (na primer, prva
godina Informatike ima šest grupa vežbi umesto dve). Više grupa znači i
više časova vežbi.

Pojam kohorte iz odeljka \ref{sec:cohorts} radi u suprotnom smeru:
dodele oblika \texttt{"groupIds": ["1i1a", "1i1b", "1i1v"]} spajaju više
grupa u jedan zajednički čas, pa nastavnik isto predavanje ne drži više
puta. Na najvećoj skali 139 od 839 časova su takve zajedničke sesije.
Zbog toga broj časova raste 16\% iako broj grupa raste 47\% -- spajanje
predavanja pojede veći deo razlike.

\subsection{Ponašanje funkcije cilja kroz vreme}

U obe konfiguracije prvo pronađeno rešenje je daleko od optimuma, i u
obe rešavač cilj spusti do nule. Bez osoblja početna vrednost je 1355,
1009 i 2440 na skalama S, M i L, a do nule se stiže kroz 23, 15 i 23
pronađena rešenja; sa osobljem početne vrednosti su 585, 1849 i 2973, a
poboljšanja ima 10, 38 i 56.

Sastav početne kazne se pri tome menja sa konfiguracijom. Bez osoblja u
njoj su samo procepi u rasporedu grupa i studenti bez mesta u učionici;
sa osobljem se na to dodaju i nastavnici koji rade više dana nego što im
je dozvoljeno, pa je početna kazna na srednjoj i najvećoj skali osetno
veća. Na najmanjoj skali je obrnuto (585 naspram 1355), što pokazuje da
te dve konfiguracije nisu ni ista instanca: sa osobljem se raspoređuje
216 časova podeljenih na 11 grupa, bez osoblja 204 časa na 8 grupa.

Poboljšanja su gusta i neravnomerna -- najveći skokovi dešavaju se rano
(na primer, sa 2259 na 1823 u svega 0,4 sekunde na skali L sa osobljem),
dok poslednji koraci do nule traju duže jer je preostale prekršaje sve
teže ukloniti: poslednjih 65 poena cilja na najvećoj skali oduzima
gotovo pola minuta pretrage.

\subsection{Skalabilnost}

Vreme do prvog dopustivog rešenja raste sa veličinom instance
(1,59 s, 4,86 s i 23,92 s bez osoblja, odnosno 2,69 s, 9,06 s i 30,03 s
sa osobljem), što je očekivano jer raste i broj sesija koje treba
smestiti u isti broj termina i učionica. Ukupno vreme do dokazano
optimalnog rešenja u konfiguraciji sa osobljem iznosi 4,05 s, 16,95 s i
121,76 s -- rast je nadlinearan, pa se vidi da najveća skala već traži
znatan deo posla u fazi optimizacije, a ne samo u traženju prvog
rešenja.

Bitno je da su svi rezultati dobijeni sa statusom \textsc{optimal},
dakle donja granica se poklopila sa nađenim rešenjem. Vremenski limit od
600 sekundi ni na jednoj skali nije ni približno iskorišćen: najsporije
pokretanje završeno je za oko dva minuta, dakle unutar 21\% raspoloživog
vremena. Prošireni model, dakle, ne samo da ostaje rešiv posle svih
dodatih ograničenja, nego i dalje ima rezervu za buduća proširenja
navedena u glavi \ref{chp:zakljucak}.

\subsection{Uticaj oskudice resursa}

Poređenje dva okruženja pokazuje da težina ovog problema mnogo više
zavisi od količine raspoloživih termina nego od broja pravila. Skraćenje
radnog dana sa dvanaest na sedam časova ne menja nijedno ograničenje
modela, a vreme rešavanja najveće instance sa osobljem raste sa 122 na
210 sekundi, dok cilj prestaje da bude dostižan u nuli.

Najviše se, međutim, produži traženje \emph{prvog} dopustivog rešenja:
na najvećoj skali sa osobljem sa 30 na 67 sekundi, a bez osoblja sa 24
na 43 sekunde. To je i suština razlike između dva okruženja -- u realnom
je težak deo posla optimizacija, a u zaoštrenom sama dopustivost. Na
manjim skalama je efekat blaži (FINAL-S sa osobljem: 2,69 s naspram
1,63 s), jer prva godina i u kratkom danu ima dovoljno prostora.

Vredi primetiti i da je početna kazna u zaoštrenom okruženju
\emph{manja} nego u realnom (2456 naspram 2973 na najvećoj skali), iako
je okruženje strože. U kratkom danu jednostavno ima manje prilika za
procepe. Razlika je u tome što se ta manja kazna ne može svesti na nulu:
u realnom okruženju cilj je velik ali potpuno uklonjiv, a u zaoštrenom
je manji ali sadrži nesvodiv ostatak od 3 poena na najmanjoj i 6 na
srednjoj i najvećoj skali.

Konačno, granica rešivosti je bliža nego što se iz ovih brojeva vidi.
Dan od 8 do 15 časova daje 1015 termina za 839 časova, dakle 83\%
popunjenosti, i rešavač tu još dokazuje optimalnost. Jedan čas kraći dan
(870 termina, 96\% popunjenosti) menja sliku iz temelja: rešavač za 600
sekundi ne nađe nijedno dopustivo rešenje. Prelaz od dokazano optimalnog
rešenja do nemogućnosti da se nađe bilo kakvo mera se, dakle, jednim
časom nastave dnevno.

\subsection{Validnost}

Na svih dvanaest pokretanja, u oba okruženja, nezavisna provera potvrđuje
da su tvrda ograničenja zadovoljena: nema dva časa u istoj učionici u
istom terminu, nijedna grupa nema dva časa istovremeno, časovi koji
zahtevaju računare su u odgovarajućim učionicama, a nijedan nastavnik
nije dvostruko zauzet.


% ------------------------------------------------------------------------------
\chapter{Zaključak i dalji rad}
\label{chp:zakljucak}
% ------------------------------------------------------------------------------

\section{Zaključak}

U ovom radu smo razmatrali problem automatskog generisanja rasporeda
časova na primeru Matematičkog fakulteta Univerziteta u Beogradu.
Implementirana su dva pristupa: programiranje ograničenja (CP-SAT) i
binarni celobrojni linearni model (0-1 ILP), poseban slučaj mešovitog
celobrojnog programiranja, rešavan SCIP-om (oznaka MIP/SCIP). Oba su
izvedena korišćenjem biblioteke Google OR-Tools.

Eksperimentalna evaluacija na tri skale realnih podataka pokazala je
sledeće:

\begin{enumerate}
\item \textbf{CP-SAT je značajno kompaktniji.}
  Sa svega 5 celobrojnih promenljivih po sesiji naspram hiljada
  binarnih promenljivih u MIP-u, CP model zauzima znatno manje
  memorije i brže se konstruiše.

\item \textbf{CP-SAT je brži.}
  Zahvaljujući globalnim ograničenjima sa specijalizovanom
  propagacijom i tehnikama učenja konflikata iz SAT rešavača,
  CP-SAT nalazi dopustivo rešenje za redove veličine brže od
  MIP/SCIP rešavača.

\item \textbf{CP-SAT bolje skalira.}
  Na većim instancama problema, prednost CP-SAT pristupa postaje
  izraženija, dok MIP/SCIP može da ne nađe rešenje u razumnom
  vremenskom roku.

\item \textbf{Oba rešavača daju validna rešenja.}
  Nezavisna validacija potvrđuje da oba pristupa, kada nađu
  rešenje, zadovoljavaju sva tvrda ograničenja.
\end{enumerate}

Na osnovu ovih rezultata, CP-SAT je izabran kao osnova za dalju
nadogradnju sistema za raspoređivanje časova. Nadogradnja je sprovedena
u dva koraka. U glavi \ref{chp:nadogradnja} uvedena je arhitektura
pravila (eng. \emph{plugin}) sa podrškom za tvrda i meka ograničenja i
funkcija cilja koja minimizuje težinsku sumu prekršaja, zajedno sa
četiri nova ograničenja kvaliteta -- spajanje časova u blokove, jedna
lokacija po danu za grupu, izbegavanje procepa u rasporedu i kapacitet
učionice kao meko ograničenje. U glavi
\ref{chp:osoblje} model je proširen nastavnim osobljem: uveden je pojam
kohorte i zajedničkih sesija (kojima se izbegava da nastavnik drži isti
čas više puta), tvrdo ograničenje da nastavnik nije dvostruko zauzet, i
tri pravila osoblja -- jedna lokacija po danu, najveći broj radnih dana i
nedeljni budžet procepa. Time je pokazano da se kompaktni CP-SAT model lako
i modularno proširuje ka realnim potrebama Matematičkog fakulteta, uz
kapacitet učionica i ograničenja nastavnika.

Završna evaluacija iz glave \ref{chp:finalna-evaluacija} pokazala je da
prošireni model ostaje praktično upotrebljiv. Iako raspoređivanje
nastavnog osoblja otprilike udvostručuje veličinu modela, na najvećoj
instanci (sve četiri godine studija, 839 časova) rešavač nalazi
dokazano optimalan raspored -- bez ijednog procepa u rasporedu grupa i
bez ijednog nastavnika preko dozvoljenog broja radnih dana -- za oko 31
sekundu, dakle unutar svega nekoliko procenata zadatog limita od 600
sekundi. Praćenje funkcije cilja kroz vreme pokazuje i da se najveći
deo kvaliteta rasporeda dobije u prvih nekoliko sekundi nakon prvog
dopustivog rešenja, što ostavlja prostora za dalja proširenja modela.

Merenje u zaoštrenom okruženju, u kome je nastava skraćena na šest
časova dnevno, pokazalo je da težina problema mnogo više zavisi od
oskudice termina nego od broja pravila: isti model tada do najboljeg
rešenja stiže za 174 sekunde umesto za 31, a najbolji cilj više nije
nula nego dokazano najmanja moguća vrednost 6. Zona u kojoj je problem
istovremeno rešiv i netrivijalan pritom je uska -- pri smanjenju broja
termina za još 17\% rešavač ne nalazi nijedno dopustivo rešenje ni za
600 sekundi.

\section{Dalji rad}

Kroz glave \ref{chp:nadogradnja} i \ref{chp:osoblje} implementirani su
meka ograničenja sa funkcijom cilja, lokacijska ograničenja za grupe i
nastavnike, grupisanje časova u blokove, raspoređivanje nastavnog osoblja
i kapacitet učionica kao meko ograničenje. Preostali pravci daljeg rada
su:

\begin{enumerate}
\item \textbf{Časovi različitog trajanja.}
  Model je ograničen na jednosatne sesije (odeljak
  \ref{sec:obim_modela}). Podrška za čas koji traje dva ili tri sata
  kao jednu jedinicu zahtevala bi da promenljiva termina nosi i
  trajanje, a da se zauzeće učionice i grupe izrazi ograničenjem nad
  intervalima, čime bi otpala potreba za pravilom o spajanju časova u
  blokove.

\item \textbf{Raspoloživost i preferencije nastavnika.}
  Uključivanje termina u kojima nastavnik ne može ili ne želi da drži
  nastavu, kao i mekih preferencija (na primer, prepodnevni termini).

\item \textbf{Dodatna meka ograničenja kvaliteta.}
  Na primer, ravnomerna raspodela časova kroz nedelju i izbegavanje
  kasnih termina, uz fino podešavanje težina u funkciji cilja i analizu
  njihovog međusobnog odnosa.

\item \textbf{Veb interfejs.}
  Izgradnja veb aplikacije koja omogućava unos podataka,
  pokretanje rešavača i vizuelni prikaz rasporeda, korišćenjem
  FastAPI okruženja za serversku stranu.
\end{enumerate}


% ------------------------------------------------------------------------------
% Literatura
% ------------------------------------------------------------------------------
\literatura

% ==============================================================================
\backmatter
% ==============================================================================

% ------------------------------------------------------------------------------
\begin{biografija}
\textbf{Dušan Trtica} rođen je 1983. godine u Beogradu. Završio je
Devetu beogradsku gimnaziju 2002. godine. Matematički fakultet
Univerziteta u Beogradu, smer informatika i računarstvo, završio je
2009. godine. Oblasti interesovanja uključuju razvoj softvera,
funkcionalno programiranje, kao i algoritme veštačke inteligencije.
\end{biografija}
% ------------------------------------------------------------------------------

\end{document}
